% ======================================================================
%  32_flesh_hecke.tex
%  骨血层之三：'t Hooft 算子 = 几何 Hecke 算子
%  主链 S7；节点 N9（§9）、N10（§10）；miracle M3
%  底本：§6.2（印刷 pp.81–89 → p082–p090）、§9（pp.123–147）、§10（pp.148–183）
%  只写本文件；不改 preamble / main / 他人文件
% ======================================================================

\section{骨血层之三：\texorpdfstring{'t Hooft}{'t Hooft} 算子就是几何 Hecke 算子}
\label{sec:flesh-hecke}

\begin{kcard}{本节的位置：主链 S7，节点 N9--N10，miracle M3}
\anchor{§9 开篇 p.123；§10 开篇 p.148；§6.2 p.81}
\textbf{一句话。}在拓扑扭化参数取 $\Psi=0$（$t=1$、$\theta=0$）时，把 $\mathcal{N}=4$ 超对称
Yang--Mills 理论紧化到 $M=\RR\times I\times C$，在 $W=I\times C$ 中一点 $p$ 处插入静态
't Hooft{} 算子 $T(\Ldual R)$。越过 $p$ 所在的时间切片 $y=y_0$ 时，膜上的全纯 $G$-丛
$E_y\to C$ 经历的跳变，\emph{恰好}是型 $\Ldual R$ 的 Hecke 修改。这就是 \textbf{S7}，
其内容一句话概括为 \textbf{M3}：“'t Hooft{} 算子作用在规范丛上的效果恰好是
Hecke 修改”。\pA\anchor{src p.127}

\textbf{本节的输入与输出。}
输入是 S5（线算子 $=$ 膜算子，§6）与 S4（$\MH$ 上的膜范畴，§5）；
输出是 S8（§11--§12 用典范余各向同性 A-膜把 Hecke 算子翻译成 $\Dmod$-模）所必需的
\emph{算子侧}词典：膜上的算子不是抽象代数生成元，而是 “几何对象的局部修改”。\pB\anchor{§8 的入口}

\textbf{为什么这段最难。}因为它把三种互不相同的对偶叠在同一处：
S-对偶（$\tau\mapsto-1/\tau$，交换 Wilson 与 't Hooft）、
GNO 对偶（$G\leftrightarrow\Ldual G$，表示 $\leftrightarrow$ 余权）、
几何 Langlands 对偶（$\MH(G)$ 的 A-膜 $\leftrightarrow$ $\MH(\Ldual G)$ 的 B-膜）。
读者最常见的失败模式，是把 \textbf{M3} 当成前两种对偶中某一种的直接推论；它不是，
它是第三种对偶在\emph{算子}层面的表现。这一点我们在 §\ref{sec:flesh-hecke-claim} 用一整段澄清。
\end{kcard}

\subsection*{符号约定}

本文件沿用 \texttt{02\_符号与风格卡} 的固定符号（$G$、$\Ldual$、$\gC$、$C$、$g$、$\MH$、$\tau$、
$t$、$\Psi$、$\phi$、$A$、$F$、$E$、$Q$、$Z(B)$ 等），只新增下表所列者。凡新增符号在
首次出现处重复一次定义；引用骨架节点一律写 $S7$、$N9$、$M3$ 这样的形式。

\begin{center}\small
\begin{tabular}{@{}lll@{}}
\toprule
\tabhead{符号} & \tabhead{含义} & \tabhead{出处}\\
\midrule
$x$，$p_0$ & $C$ 上的插入点（$C$ 为紧 Riemann 面，亏格 $g$） & \anchor{src p.128}\\
$W=I\times C$ & 三维流形；$I$ 为区间，$C$ 为 Riemann 面 & \anchor{src p.149}\\
$y$ & $I$ 的坐标；$E_y=E|_{\{y\}\times C}$ & \anchor{src p.129}\\
$\mathcal{O}(x)$ & $C$ 上在 $x$ 外全纯、至多 $x$ 处单极点的函数对应线丛 & \anchor{src p.130}\\
$\mathcal{Y}(\Ldual R)$，$\mathcal{Y}(\Ldual w;x)$ & 型 $\Ldual R$（或最高权 $\Ldual w$）的 Hecke 修改空间 & \anchor{src p.137}\\
$\mathrm{Gr}_G$ & 仿射 Grassmann 流形 $G_{\CC}((z))/G_{\CC}[[z]]$ & \anchor{src p.146}\\
$\mathcal{Z}(\Ldual w_1,p_1;\dots)$ & 带指定奇异性 Bogomolny 方程解的模空间 & \anchor{src p.156}\\
$\mathrm{IH}^\bullet$ & 交集上同调 & \anchor{src p.174}\\
$K$ & A-模型的幽灵数（ghost number） & \anchor{src p.173}\\
$N_\alpha$ & 张量积重数空间 $\mathrm{Hom}_{\Ldual G}(\Ldual R\otimes\Ldual R',\Ldual R_\alpha)$ & \anchor{src p.167}\\
$\mathcal{Q}$ & Higgs 丛的 Hecke 对应（不用 $Q$，避免与超荷冲突） & \anchor{src p.182}\\
$\rho$ & 同态 $U(1)\to G$；其共轭类即 GNO 荷 & \anchor{src p.84}\\
$\Ldual w=(m_1,\dots,m_N)$ & $\Ldual G$ 的最高权（对 $G=U(N)$ 为整数 $N$-plet） & \anchor{src p.131}\\
$\Delta$ & 修改空间的复维数 & \anchor{src p.144}\\
\bottomrule
\end{tabular}
\end{center}

\subsection{导言：三层任务与一条主线}\label{sec:flesh-hecke-intro}

这一段（原文 §6.2、§9、§10）要建立的东西可以拆成三层，而这三层在原文里是
\emph{交织}着写的：§6.2 给物理定义，§9 给数学对象并直接宣布二者相等，
§10 才回过头用 Bogomolny 方程把“相等”变成模空间的相等。我们的任务是把它拆开、
把被跳过的步骤补上、并标出哪些地方原文实际上是“用物理论证代替了证明”。

\begin{keyidea}
\textbf{三层任务。}
\begin{enumerate}
\item \textbf{物理层}（§6.2）：'t Hooft{} 算子是\emph{无序算子}（disorder operator），
  它不测量场的某个值，而是规定场沿一条线的奇异性；插入了它的路径积分，是在
  \emph{带奇异性}的场上做的。\pA\anchor{src p.81}
\item \textbf{数学层}（Beilinson--Drinfeld）：Hecke 修改是把 $C$ 上一点 $x$ 处的
  $G$-丛 $E$ 换成 $E'$，使 $E'$ 与 $E$ 在 $C\setminus\{x\}$ 上同构，而 $x$ 处的
  “差”由一个仿射 Grassmann 流形的 Schubert 胞腔的参数刻画。\pA\anchor{src p.146}
\item \textbf{相等}（§9.1、§10.3）：4d 理论紧化到 2d 后，'t Hooft{} 算子对膜上的
  全纯丛的作用，就是第二层的 Hecke 修改；并且这个相等可以被提升为
  \emph{模空间的相等}：Bogomolny 方程解空间（带边界条件）$\mathcal{Z}$ 与
  修改空间 $\mathcal{Y}$ 是同一个复流形。\pA\anchor{src p.162}
\end{enumerate}
\end{keyidea}

\textbf{为什么必须拆成三层。}
如果只读 §9.1，读者会觉得 \textbf{M3} 是一个\emph{奇迹}：物理上插入一根奇异的线，数学上丛就
被修改了。但 \textbf{M3} 不是奇迹，它是\emph{可以逐条验证}的：'t Hooft{} 算子的奇异数据是一个
同态 $\rho:U(1)\to G$；Hecke 修改的类型是仿射 Grassmann 中一个 $G_{\CC}([[z]])$-轨道；
而这两者是\emph{同一个} $\rho$（其一为实数形式，另一为其复化）。这个“同一个 $\rho$”
就是全篇的枢纽，我们在 §\ref{sec:flesh-hecke-gr} 把它写成一条显式的等式。\pB\anchor{src p.147}

\textbf{一条主线（建议读者按此顺序读）。}
\begin{enumerate}
\item §\ref{sec:flesh-hecke-thooft}：先弄清 't Hooft{} 算子是什么，尤其是它\emph{为什么不是}
  一个通常意义上的算子（这一点原文 p.123 说得很明确，却常被忽略）。
\item §\ref{sec:flesh-hecke-hecke}：Hecke 修改的定义与两个最小例子
  （$G=GL_1$、$G=GL_2$），以及在什么意义上它是“局部”的。
\item §\ref{sec:flesh-hecke-claim}：把 §9.1 的论证链逐条重建成七个命题，
  各标 $A$/$B$，得到 M3 的确切陈述。
\item §\ref{sec:flesh-hecke-gr}：仿射 Grassmann 与 Schubert 分层，说明修改类型的
  参数空间是什么。
\item §\ref{sec:flesh-hecke-s10}：§10 的四件技术工作（边界条件、单极子冒泡、
  Kähler 商、OPE 与扩展 Bogomolny 方程），以及
\item §\ref{sec:flesh-hecke-honest}：怎样诚实地读这一段——哪些是定理，哪些是物理论证。
\end{enumerate}

\textbf{五个最容易犯的错。}（'t Hooft{} 不是 Wilson 的对偶那么简单；Hecke 不是换掉丛；
不是所有修改都等于张量 $\mathcal{O}(x)$；“算子”不按字面读；M3 不是前两种对偶的推论。）
这里不展开——后文给了五张“排除错误直觉”卡片，逐条对照即可。


\subsection{第一层：\texorpdfstring{'t Hooft}{'t Hooft} 算子的物理定义（§6.2，pp.~81--89）}
\label{sec:flesh-hecke-thooft}

\subsubsection{动机：为什么必须发明一种新的算子}

在 $\Psi=\infty$（$t=\pm\ii$）的拓扑扭化理论里，Wilson 算子是现成的：取 $G$ 的一个表示 $R$、
一条闭曲线 $S$，算 Wilson 算子就是算 $S$ 上联络的和乐（holonomy）在 $R$ 中的迹。它是
\term{order operator}{有序算子}：先有一个经典表达式（和乐），再把它当量子算子用。\pA\anchor{src p.81}

但 S-对偶要求有对偶对象。对偶变换 $S:\Psi\mapsto-1/n_g\Psi$ 把 $\Psi=\infty$ 送到
$\Psi=0$；于是“$\Psi=0$ 处的理论里的某个算子”应当与 Wilson 对偶。问题在于：
$\Psi=0$ 时理论里没有任何现成的和乐型表达式可用——扭化参数取这个值时，被保留的
超荷是另一个，能写成经典表达式的量变了。\pA\anchor{src p.81}

出路是\emph{反过来}定义算子：不规定“沿 $S$ 的和乐等于什么”，而规定
“场在 $S$ 附近必须怎样奇异”。这就是 \term{disorder operator}{无序算子}：
它的定义不是作用的表达式，而是路径积分的一个\emph{约束}（边界条件）；
然后\emph{在满足该奇异性的场空间中}做路径积分。原文的原话是：无序算子的定义方式
是“规定经典场应当具有的奇异性，并在奇异性存在的情况下做路径积分”。
\pA\anchor{src p.81}

\begin{keyidea}
\textbf{两类算子的对偶是“同一件事的两种描述”，不是“两个算子相乘”。}
在量子力学里，有序算子与无序算子可以在适当情形下有相同的代数性质，并且常常在对偶下
互换。Wilson 的对偶（在 $\Psi=0$ 的理论里）就是 't Hooft{} 算子。
这句话的\emph{内容}是：$\Psi=\infty$ 一侧的表示论数据（$G$ 的表示 $R$）
与 $\Psi=0$ 一侧的表示论数据（$\Ldual G$ 的表示 $\Ldual R$）互为对偶，
不是“同一个算子的两个名字”。\pA\anchor{src p.81}
\end{keyidea}

这里有一个容易被忽略的技术细节，它在本节的公式里反复出现：定义 Wilson 算子需要
曲线 $S$ 的\emph{定向}；而定义 't Hooft{} 算子需要 $S$ 的\emph{法丛}的定向。
把 $S$ 的小邻域等同于 $S\times\RR^3$、记 $\epsilon_3$ 为 $\RR^3$ 的体积形式、$r$ 为到
$S$ 的距离，则
\begin{equation}
\epsilon_3=\dd r\,\dd\mathrm{Vol},
\label{eq:eps3}
\end{equation}
't Hooft{} 算子的定义只依赖法丛的定向。\pA\anchor{src p.81}\anchor{eq 6.6}

\subsubsection{最小例子：$G=U(1)$ 与奇异 Dirac 单极子}

先看最简单的：\term{Dirac monopole}{Dirac 单极子}。$\RR^4=\RR\times\RR^3$，坐标
$(s,\vec x)$，奇异沿直线 $L=\{\vec x=0\}$。$\mathrm{U}(1)$ 的 Maxwell 方程是
$\dd F=\dd\star F=0$，它在 $\RR^3\setminus\{0\}$ 上有旋转与平移对称的解
\begin{equation}
F=\frac{\ii}{2}\star\dd\frac{1}{|\vec x|}.
\label{eq:dirac}
\end{equation}
（$\star$ 表明这个解依赖法丛的定向。）把它沿 $s$ 方向拉回到 $\RR^4\setminus L$ 即可。
归一化选得使包围奇点的 $S^2$ 上
\begin{equation}
\oint\frac{\ii F}{2\pi}=1
\qquad\Longleftrightarrow\qquad
\oint_{Y}c_1(L)=1,
\label{eq:flux1}
\end{equation}
即存在一个曲率为 $F$ 的酉线丛 $L\to\RR^4\setminus L$，其第一陈数在包围奇点的球面上为 $1$。\pA\anchor{src p.83}
\anchor{eq 6.7}

于是 $G=U(1)$ 的基本 't Hooft{} 算子 $T(0)$ 定义为：\emph{在 $M\setminus S$ 上对
abelian 规范场做路径积分，并要求靠近 $S$ 时场与奇异 Dirac 单极子一致}；
等价的微分条件是 $F$ 在 $S$ 附近发散为
\begin{equation}
F\sim\frac{\ii}{2}\star_3\dd\frac{1}{r},
\label{eq:sing}
\end{equation}
其中 $\star_3$ 是法向平面上的 Hodge 星。更一般地取整数 $m$，要求奇异部分为
$(\ii m/2)\star_3\dd(1/r)$，就得到磁荷 $m$ 的算子 $T(m)$。\pA\anchor{src p.82}
\anchor{src p.83}

这里已经埋下了本节最重要的伏笔：'t Hooft{} 算子的\emph{全部}内容就是
“把路径积分的定义域从\{全体场\}限制到\{在 $S$ 附近有指定奇异性的场\}”。
它没有“作用在态上”的公式，而是一个\emph{几何约束}。后面我们会看到，
在紧化到二维之后，这个几何约束恰好可以用“丛的局部修改”来翻译。

原文还补了一句拓扑上的警告：若 $S$ 的同调类在 $H_1(M,\ZZ)$ 中非零，则 $\mathrm{U}(1)$
理论里根本不存在这样的规范场，路径积分取零；对偶地，Wilson 算子的期望值也为零。
这说明 't Hooft{} 算子的“定义域”本身就带着拓扑信息。\pA\anchor{src p.82}

\subsubsection{拓扑对称要求奇异的是整个超多重态，不只是规范场}

上面的定义只规定了 $F$ 的奇异性。但 $\Psi=0$ 处我们要求算子保持\emph{拓扑对称}
（即扭化后保留下来的那个标量超荷）。判据是：'t Hooft{} 算子保持拓扑对称，当且仅当
场在 $S$ 附近的奇异行为与拓扑对称相容。对 abelian 情形取 $t=1$、$\tau$ 虚，
保持拓扑对称的场条件是
\begin{equation}
(F+\dd\phi)_+=(F-\dd\phi)_-=0,
\label{eq:topysym}
\end{equation}
而若 $F$ 有 Dirac 奇异、$\phi$ 光滑，这个条件不满足。\pA\anchor{src p.83}\anchor{eq 6.8}

因此 $\phi$ 也必须奇异。取 $s$ 为沿 $L$ 的“时间”坐标，
\begin{equation}
F=\frac{\ii}{2}\star_3\dd\frac{1}{|\vec x|},
\qquad
\phi=\frac{\ii}{2|\vec x|}\,\dd s,
\label{eq:phising}
\end{equation}
则 \eqref{eq:topysym} 被满足。\pA\anchor{src p.83}\anchor{eq 6.9}

这一步是很多读者会漏掉的：'t Hooft{} 算子不是“只对规范场设条件”，
它必须把 Higgs 场（或者说 $\phi$ 的相应分量）也一起做成奇异的，否则它破坏拓扑对称、
在拓扑理论里根本不合法。这个“$\phi$ 也要奇异”的事实，正是 §\ref{sec:flesh-hecke-s10}
里扩展 Bogomolny 方程（$\phi$ 有全纯截面结构）的物理起源。\pB\anchor{src p.83}

\subsubsection{一般 $G$：同态 $\rho:U(1)\to G$ 与 GNO 荷}

推广到任意紧致李群 $G$：取任意同态 $\rho:U(1)\to G$，要求场沿 $S$ 的奇异性是
\eqref{eq:phising} 在 $\rho$ 下的像。算子只依赖 $\rho$ 的共轭类。\pA\anchor{src p.84}

关键的表示论事实是：$\rho$ 的共轭类由 $\Ldual G$ 的最高权分类，也就是由 $\Ldual G$
的不可约表示分类。因此对每个不可约表示 $\Ldual R$ 与每条一维曲线 $S$，有一个拓扑
't Hooft{} 算子 $T(\Ldual R,S)$，简写 $T(\Ldual R)$。\pA\anchor{src p.84}

原文把这句话称为“经典 GNO 对偶”（Goddard--Nuyts--Olive），并指出参考文献
[90] 的贡献在于把它\emph{重述为算子}而不是态：GNO 原本分类的是磁单极子的荷，
而这里同一个组合数据被用来定义一类线算子。\pA\anchor{src p.84}

对 $G=U(N)$，同态 $\rho$ 可以取成
\begin{equation}
\rho(\ee^{\ii\alpha})=\mathrm{diag}\bigl(\ee^{\ii m_1\alpha},\dots,\ee^{\ii m_N\alpha}\bigr),
\qquad m_1\ge m_2\ge\cdots\ge m_N,
\label{eq:rhoUN}
\end{equation}
于是 't Hooft{} 算子由整数 $N$-plet $\Ldual w=(m_1,\dots,m_N)$（模置换）刻画；
在 Lie 代数层面 $\rho$ 把 $1\in\mathfrak{u}(1)$ 送到 $\mathrm{diag}(m_1,\dots,m_N)$。\pA\anchor{src p.131}
\anchor{eq 9.11}
相应地，算子 $T(\Ldual w)$ 的奇异条件是：$p$ 附近曲率的奇异部分为
\begin{equation}
F\sim\frac{\ii}{2}\star\dd\frac{1}{|\vec x-p|}\cdot\mathrm{diag}(m_1,\dots,m_N).
\label{eq:Fdiag}
\end{equation}
\pA\anchor{src p.132}\anchor{eq 9.12}

\subsubsection{IR 拓扑版本为什么可以被理解为“对规范丛的修改”}

这是任务要求的第一问，也是本小节最需要重建的一步。原文在 p.84 只用了一句话把它带过：
't Hooft{} 算子的分类“与 Grothendieck 的 $\mathbb{P}^1$ 上 $G$-丛分类密切相关”。
我们把它展开成三步。

\begin{keyidea}
\textbf{三步论证。}
(i) 把 $S^2$ 等同于包围奇异点的球面：构造 $T(\Ldual R)$ 用的 abelian 解在 $S^2$ 上给出
一个 $G$-丛，而 $\mathbb{P}^1$ 上全纯 $G$-丛正由 $\mathcal{O}(1)$ 经 $\rho:U(1)\to G$
（复化后 $\rho_{\CC}:\CC^*\to G_{\CC}$）给出。\pA\anchor{src p.84}
(ii) 反向也成立：$\mathbb{P}^1$ 上每个 $G$-丛（即使不稳定）都有 Yang--Mills 联络，
而每个 YM 解都等价于构造 't Hooft{} 算子用的 abelian 解。\pA\anchor{src p.84}
(iii) 故 't Hooft{} 算子的分类 $=$ $\mathbb{P}^1$ 上 $G$-丛的分类；同一句话在四维是
“在 $S$ 附近保留一个指定的丛”，在二维是“在点 $x$ 处保留一个指定的丛”。
\pA\anchor{src p.84}
\end{keyidea}

第三步的“二维表述”是本文的关键：一旦理论紧化到 $C$ 上（$M=\RR\times I\times C$），
包围奇异线的球面 $S^2$ 就被换成了 $C$ 上包围点 $x$ 的\emph{小圆盘}。于是
“$\Psi=0$ 时保持拓扑对称的 't Hooft{} 算子”在二维图像里就是：
“在 $x$ 的小圆盘上换一个丛，圆盘外不动”。这正是 Hecke 修改的定义。
这个翻译是 $B$ 级重建：原文没有把它写成一句命题，但 p.84 的“分类相同”加上
§9.1 的具体计算足以支撑它。\pB\anchor{src p.84；重建依据：p.132 的 $\oplus L_i$ 分解与 p.146 的
仿射 Grassmann 表述}

还有一个拓扑细节值得单独记下：$\mathbb{P}^1$ 上的全纯 $G$-丛可以拓扑非平凡。
$G$-丛的拓扑由 $\xi\in H^2(S^2,\pi_1(G))\cong\pi_1(G)$ 分类，而
$\pi_1(G)\cong Z(\Ldual G)$；构造 $T(\Ldual R)$ 所用的球面上的丛非平凡，
当且仅当 $\Ldual G$ 的中心在 $\Ldual R$ 上非平凡作用。\pA\anchor{src p.84}
用 \eqref{eq:rhoUN} 的语言说：$\sum_i m_i\ne 0$ 时 $T(\Ldual w)$ 会改变 $G$-丛的拓扑
（因为 $c_1\mapsto c_1+\sum_i m_i$，见 \eqref{eq:c1shift}）；$\sum_i m_i=0$ 时拓扑不变。
这个判据在 §\ref{sec:flesh-hecke-claim} 会以“分解后每根线丛各自被张量化”的形式
再次出现。\pB\anchor{src p.133}

\begin{kcard}{排除错误直觉 1：“'t Hooft{} 算子就是 Wilson 算子的对偶”这句话不能这样用}
\anchor{src p.81；src p.87}
\textbf{错误理解：}既然 S-对偶交换 $\Psi=\infty$ 与 $\Psi=0$，那么把 Wilson 算子
（和乐的迹）换成“对偶表示的和乐”就得到 't Hooft{} 算子。

\textbf{为什么错：}
(i) 't Hooft{} 算子的定义不是任何经典量的量子化，而是路径积分的\emph{定义域约束}；
在 $\Psi=0$ 的理论里没有“和乐型”的表达式可以替换。\pA\anchor{src p.81}
(ii) 只有\emph{有理} $\Psi$ 时才能把 Wilson 与 't Hooft{} 拼成一个保持拓扑对称的
复合算子 $WT_{n,m}$；一般 $\Psi$ 下 't Hooft{} 算子破坏拓扑对称，需要乘一个
Wilson 因子才能修复，而规范不变性又要求 $n=-m\Psi$ 为整数。所以
“纯 't Hooft{} 算子”只在 $\Psi=0$ 这一个点上作为拓扑算子存在。\pA\anchor{src p.87}
\anchor{src p.88}
(iii) 更根本的是：S-对偶只是把一侧的算子映到另一侧的算子；而 \textbf{M3}（'t Hooft $=$ Hecke）
说的是“算子作用于膜”的效果，属于另一个对偶（几何 Langlands）。见§\ref{sec:flesh-hecke-claim}。

\textbf{正确图像：}'t Hooft{} 算子是一个\emph{边界条件}。它的“对偶性”体现在
两个侧面：作为奇异性它由 $\Ldual G$ 的表示分类（GNO 侧）；作为二维图像它由
仿射 Grassmann 的 Schubert 胞腔分类（Hecke 侧）。这两侧是同一枚硬币。
\end{kcard}

\paragraph{自检问题。}
(1) 为什么 Wilson 算子需要 $S$ 的定向，而 't Hooft{} 算子需要法丛的定向？
提示：看 \eqref{eq:dirac} 里的 $\star$。(2) 若 $S$ 的同调类非零，为什么 $\mathrm{U}(1)$
理论里 't Hooft{} 算子的路径积分取零，而 Wilson 的期望值也为零？
(3) 在 \eqref{eq:rhoUN} 中若把 $m_i$ 全加上同一个整数 $c$，算子变成什么？
（答案在 §\ref{sec:flesh-hecke-hecke}：相当于把输出丛整体张量 $\mathcal{O}(x)^c$。）


\subsection{第二层：Hecke 修改的数学定义（Beilinson--Drinfeld）}
\label{sec:flesh-hecke-hecke}

\subsubsection{定义：一对丛，加上一个“在挖去一点后成立”的同构}

设 $C$ 是复光滑射影曲线（本文只用到 $C=\PP^1$ 与一般紧 Riemann 面），$x\in C$ 是点，
$E\to C$ 是 $G$-丛（对 $G=GL_N$ 就是秩 $N$ 全纯向量丛）。所有讨论都在 $x$ 的
\emph{形式邻域}里进行：取局部坐标 $z$，形式圆盘 $\mathcal{O}=\CC[[z]]$、
形式穿孔圆盘 $\mathcal{K}=\CC((z))$。

\begin{definition}[Hecke 修改]\label{def:hecke}
$E$ 在 $x$ 处的\textbf{Hecke 修改}是一个二元组 $(E',\beta)$，其中 $E'\to C$ 是另一个
$G$-丛，$\beta:E'|_{C\setminus\{x\}}\xrightarrow{\ \sim\ }E|_{C\setminus\{x\}}$ 是
$C\setminus\{x\}$ 上的同构。若 $\beta$ 可以延拓为 $C$ 上的同构，则称该修改是
\textbf{平凡的}（此时 $E'\cong E$，与 $E$ 相同）。
\end{definition}

这个定义里最要紧的一点，是 $\beta$ 是数据的一部分：我们比较的不是“两个丛”，
而是“两个丛连同它们在 $C\setminus\{x\}$ 上的同构”。这一点决定了后面所有
“等价”与“分类”的陈述方式。

\begin{definition}[修改的类型]\label{def:type}
取 $E$ 在 $x$ 附近的一个局部平凡化（形式圆盘上的标架）。在这个标架下，$E'$ 由
转移函数 $g(z)\in G(\mathcal{K})$ 给出，而 $E'$ 在圆盘上的平凡化的改变是右乘
$G(\mathcal{O})$。另一方面，局部平凡化 $E$ 的选择带来左乘 $G(\mathcal{O})$ 的自由。
于是\textbf{修改的类型}是双陪集
\begin{equation}
G(\mathcal{O})\backslash G(\mathcal{K})/G(\mathcal{O}),
\label{eq:doublecoset}
\end{equation}
而所有可能的修改（固定 $E$、固定 $x$，模去上述两类平凡化）构成
\term{affine Grassmannian}{仿射 Grassmann 流形}
\begin{equation}
\mathrm{Gr}_G=G(\mathcal{K})/G(\mathcal{O}).
\label{eq:grdef}
\end{equation}
\end{definition}

对 $G$ 单连通、$T\subset G$ 为极大环面时，\eqref{eq:doublecoset} 的集合由支配余权
$\lambda\in X_*(T)^+$ 分类；而 $X_*(T)^+$ 与 $\Ldual G$ 的支配最高权是同一组组合数据。
\textbf{这句话就是 \textbf{M3} 的组合骨架}：'t Hooft{} 算子的 GNO 荷（$\Ldual G$ 的最高权）
与 Hecke 修改的类型（支配余权）是同一组数据。我们在
§\ref{sec:flesh-hecke-gr} 把它写成显式等式。\pA\anchor{src p.146}

\subsubsection{最小例子一：$G=GL_1$}

$G=\CC^*$、$E=\mathcal{O}_C$ 平凡线丛。设 $(E',\beta)$ 是 $x$ 处的修改。
在 $C\setminus\{x\}$ 上 $E'\cong\mathcal{O}$，故 $E'$ 是一个除子类为 $k x$
的线丛，$k\in\ZZ$ 唯一；$E'=\mathcal{O}(kx)$。于是：

\begin{proposition}[$GL_1$ 的修改]\label{prop:gl1}
$GL_1$ 在 $x$ 处的修改类型的集合是 $\ZZ$；型 $k$ 的修改唯一
（空间 $\mathcal{Y}(k)$ 是单点），且 $\deg E'=\deg E+k$。在局部坐标下，
型 $k$ 的修改就是把局部标架乘以 $z^{-k}$：
$E'$ 的截面形如
\begin{equation}
s(z)=z^{-k}f(z),\qquad f\in\mathcal{O}.
\label{eq:gl1section}
\end{equation}
\end{proposition}

\begin{proof}[推导（原文完全省略）]
$\mathrm{Gr}_{GL_1}=\CC^*(\mathcal{K})/\CC^*(\mathcal{O})$。$\CC^*(\mathcal{K})$ 中的元素是
$z^k$ 乘一个单位（$k\in\ZZ$），而单位可被 $\CC^*(\mathcal{O})$ 吸收，故
$\mathrm{Gr}_{GL_1}\cong\ZZ$，由 $z$ 的赋值给出。几何上，$E'|_{C\setminus x}\cong E|_{C\setminus x}$
给出一个在 $x$ 外无零点无极点的截面的类，它在 $x$ 处的零点阶（可正可负）就是 $k$，
于是 $E'=E\otimes\mathcal{O}(x)^k=\mathcal{O}(kx)$。
\end{proof}

注意 \eqref{eq:gl1section} 已经说明了“局部修改”的含义：型 $k$ 的修改
只改变 $x$ 附近截面的\emph{极点允许阶数}，而在 $x$ 之外 $E'$ 与 $E$ 有\emph{完全相同}的截面。
但也要立刻看到：$k\ne0$ 时 $E'\not\cong E$（次数不同），所以“局部”
绝不是说“全局对象不变”。\pB\anchor{src p.137}

\subsubsection{最小例子二：$G=GL_2$，型 $(1,0)$ 与 $(2,0)$}

设 $E=\mathcal{O}\oplus\mathcal{O}$，局部标架 $s_1,s_2$ 在 $x$（取 $x$ 为 $z=0$）附近全纯。
型 $\Ldual w=(1,0)$ 的修改（这是原文 §9.2 的第一个例子）定义为：$E'$ 由
\begin{equation}
t_1=z^{-1}s_1,\qquad t_2=s_2
\label{eq:gl2frame}
\end{equation}
生成。$E'$ 在 $x$ 附近的任意全纯截面因此形如
\begin{equation}
s(z)=\frac{c}{z}\,s_1(0)+\bigl(h_1(z),h_2(z)\bigr),
\label{eq:gl2section}
\end{equation}
其中 $c\in\CC$、$h_1,h_2$ 在 $0$ 附近全纯。\pA\anchor{src p.137}\anchor{eq 9.16}

\begin{proof}[推导：为什么 \eqref{eq:gl2section} 里只剩下 $s_1(0)$]
$E'$ 的截面是 $g_1(z)t_1+g_2(z)t_2$（$g_i$ 全纯），即
$g_1(z)z^{-1}s_1(z)+g_2(z)s_2(z)$。把 $s_1(z)=s_1(0)+zs_1'(0)+O(z^2)$ 代入：
\begin{equation*}
g_1z^{-1}s_1
=\frac{g_1(0)}{z}s_1(0)+g_1'(0)s_1(0)+g_1(0)s_1'(0)+O(z).
\end{equation*}
后两项全纯，可以并进 $h_1$；所以只有 $c=g_1(0)$ 这一项真正在 $x$ 处产生极点。
又 $c\to\lambda c$ 可由 $E'$ 在圆盘上平凡化的改变（乘以常数矩阵）吸收，故
\eqref{eq:gl2section} 中真正有意义的参数是 $s_1(0)$ 到复标度的类。
\end{proof}

$s_1(0)\in\CC^2\setminus\{0\}$ 模 $\CC^*$ 正是 $\PP^1$ 的齐次坐标。于是
\begin{equation}
\mathcal{Y}\bigl((1,0);x\bigr)\cong\PP^1=\mathrm{Gr}(1,2).
\label{eq:y10}
\end{equation}
这是“修改空间”的第一个不平凡例子：它不是单点，而是一个紧致光滑的射影簇，
参数是“$E'$ 相比 $E$ 在哪个方向上多了一个极点”。\pA\anchor{src p.137}

型 $(2,0)$ 的修改（$E'$ 由 $z^{-2}s_1,s_2$ 生成）给出
\begin{equation}
s(z)=\frac{c}{z^{2}}\bigl(u+zu',v+zv'\bigr)+\frac{d}{z}(u,v)+\bigl(h_1(z),h_2(z)\bigr),
\label{eq:y20}
\end{equation}
其中 $(u,v)=s_1(0)$，$(u',v')$ 来自 $s_1$ 的一阶项。两个等价关系
$(u,u',v,v')\to\lambda(u,u',v,v')$ 与 $(u',v')\to(u',v')+w(u,v)$ 把参数空间压成
带权的射影空间 $\mathbb{WP}^2(1,1,2)$ 挖去一点。
原文随后证明：挖去的那个点有自然的补法，补上之后得到的紧化
$\mathcal{Y}((2,0);x)$ 恰好也包含型 $(1,1)$ 的修改。\pA\anchor{src p.138}
\anchor{src p.139}

\begin{keyidea}
\textbf{$(2,0)$ 的紧化包含 $(1,1)$，这件事的物理意义是“不稳定性”。}
型 $(2,0)$ 的修改对应的 $\PP^1$ 上丛是 $\mathcal{O}(2)\oplus\mathcal{O}$，
它\emph{不稳定}，可以无穷小扰动成全纯型 $\mathcal{O}(1)\oplus\mathcal{O}(1)$；
后者正是型 $(1,1)$ 的修改（张量 $\mathcal{O}(x)$）。所以为了让 A-模型的上同调有话可说，
必须把修改空间紧化，而紧化就自动把“不稳定退化”的型收进边界。
型 $(1,0)$ 的丛 $\mathcal{O}(1)\oplus\mathcal{O}$ 不能这样扰动，故它的修改空间
$\PP^1$ 已经紧致光滑。\pA\anchor{src p.139}
\end{keyidea}

\subsubsection{“局部”的精确含义}

上面两个例子已经把“局部修改”的意思说清楚了，但这个词太容易被误读，值得单独钉死。
\begin{enumerate}
\item \textbf{数据集中在一点。}修改由 $(E',\beta)$ 给出，$\beta$ 定义在 $C\setminus\{x\}$ 上；
  也就是说除了 $x$，$E'$ 与 $E$ 没有任何差别（通过 $\beta$ 认同）。
\item \textbf{修改的类型是离散标签 $+$ 有限维连续参数。}型 $\lambda$（支配余权）是离散的；
  同一型内的连续参数构成一个有限维簇 $\mathcal{Y}(\lambda)$，典型地是 Schubert 簇。
  't Hooft{} 算子的路径积分，物理上就是对 $\mathcal{Y}(\lambda)$ 积分（更准确地说，
  是对以它为典型纤维的模空间积分，见 §\ref{sec:flesh-hecke-s10}）。
\item \textbf{但全局拓扑可以改变。}$c_1$ 可以变（\eqref{eq:c1shift}），稳定性可以变。
  “局部”说的是\emph{修改的自由度}是局部的，不是“拓扑不变”。
\end{enumerate}

\begin{kcard}{排除错误直觉 2：Hecke 修改不是“把丛换掉”}
\anchor{src p.136；src p.146}
\textbf{错误理解：}Hecke 修改就是“在 $x$ 处把 $E$ 换成另一个丛 $E'$”。

\textbf{为什么错：}照这个说法，$E$ 与 $E'$ 之间没有“关系”，那“修改”
就无从比较，也不可能有一个参数空间。正确的对象是二元组 $(E',\beta)$：
$\beta$ 是 $C\setminus\{x\}$ 上的同构，是分类数据的一部分。两个修改 $(E',\beta)$ 与
$(E'',\beta')$ \emph{等价}，当且仅当 $\beta'\circ\beta^{-1}$ 延拓过 $x$ 成为
$C$ 上的同构。$\beta$ 延拓不过 $x$ 的“程度”，就是修改的类型。\pA\anchor{src p.136}

\textbf{正确图像：}Hecke 修改是“一个在 $C\setminus\{x\}$ 上成立、在 $x$ 处失效的同构”。
失效的方式（用 $G(\mathcal{O})\backslash G(\mathcal{K})/G(\mathcal{O})$ 的哪个双陪集）是离散类型；
失效的连续参数（在 Schubert 胞腔里的位置）是修改空间的坐标。
\end{kcard}

\paragraph{自检问题。}
(1) 在 \eqref{eq:doublecoset} 中，为什么左右两边都要模 $G(\mathcal{O})$？
分别对应哪种平凡化选择的自由？
(2) 对 $G=GL_2$，写出型 $(1,1)$ 的修改，并验证它等于 $E\otimes\mathcal{O}(x)$。
(3) 型 $(2,0)$ 与型 $(1,1)$ 的修改空间维数分别是多少？为什么前者非紧而后者紧？


\subsection{第三层：核心论断——\texorpdfstring{'t Hooft}{'t Hooft} 算子作用在膜上就是 Hecke 修改}
\label{sec:flesh-hecke-claim}

\subsubsection{设定：把四维理论放到 $M=\RR\times I\times C$ 上}

现在执行 S3（紧化）在本题中的具体形式。取
\begin{equation}
M=\RR\times W,\qquad W=I\times C,
\label{eq:Msetup}
\end{equation}
$I$ 是区间、$C$ 是 Riemann 面。考虑\emph{静态} 't Hooft{} 算子：它的支撑是
$\RR\times\{p\}$，$p=y_0\times x$，其中 $y_0\in I$、$x\in C$。
在 $M$ 上把时间取为 $\RR$ 的坐标 $s$，扭化参数取 $\Psi=0$（即 $t=1$、$\theta=0$）。\pA\anchor{src p.123}
\anchor{src p.125}

$\Psi=0$ 的超对称条件是
\begin{equation}
(F-\phi\wedge\phi+t\dd\phi)_+=0,\qquad
(F-\phi\wedge\phi-t^{-1}\dd\phi)_-=0,\qquad
D{*}\phi=0,
\label{eq:bps}
\end{equation}
$t=1$ 时前两式合写为 $F-\phi\wedge\phi+\star\dd\phi=0$。用复联络
$\mathcal{A}=A+\ii\phi$，等价地写成 $F+\ii\star F=0$ 型（自对偶）方程。\pA\anchor{src p.124}
\anchor{eq 9.1}

原文随后做了两步简化，得到本节的主战场：
\begin{enumerate}
\item 若只考虑同时保持复结构 $J$ 与 $K$ 的 A-模型拓扑超荷的解，则解\emph{时间无关}；
  证明思路是用 $\phi_1$（$\phi$ 在 $I$ 方向的分量）的消失定理：乘 $\phi_1$ 取迹积分，
  被积式负定，故 $\phi_1=0$。\pA\anchor{src p.126}\anchor{eq 9.3}
\item 再假定解在时间反演 $T$ 下不变，则 $G$-丛与联络从 $W$ 拉回、$\phi=\phi_0\dd s$，
  方程约化为
  \begin{equation}
  F=\star\dd\phi_0,
  \label{eq:bogomolny}
  \end{equation}
  即标准 \term{Bogomolny equation}{Bogomolny 方程}（BPS 单极子方程）。\pA\anchor{src p.127}
  \anchor{eq 9.5}
\end{enumerate}

这两步是\emph{假设}而非结论：原文明确说，$\phi_1$ 的消失定理对最重要的 (B,A,A) 膜成立，
而时间反演不变性则把边界条件限制得很死；真正对几何 Langlands 有用的是后面 §10.5 的
扩展 Bogomolny 方程。\$9 先做这个简化情形，因为它已足够展示主要想法。\pA\anchor{src p.125}
\anchor{src p.179}

于是问题变成：在 $W=I\times C$ 上解 \eqref{eq:bogomolny}，在 $p=y_0\times x$ 处有
型 $\Ldual w$ 的奇异性，两端取适当的边界条件（左端指定丛 $E_-$，右端 $\phi_0=0$），
问右端的丛 $E_+$ 是什么。\pA\anchor{src p.149}

\subsubsection{§9.1 的论证链：七步重建}

原文 §9.1 的论证分成七步，其中前五步是原文明确做的计算，后两步是原文一句话带过、
需要补上“为什么”的。我们逐条列出并标注 $A$/$B$。

\paragraph{第 1 步（A）：远离 $y_0$ 时，$E_y$ 的全纯类型与 $y$ 无关。}
把 $E$ 限制到 $\{y\}\times C$ 得全纯丛 $E_y\to C$。取 $A_y=0$ 的规范，则
$F_{yz}=\partial_yA_z$；Bogomolny 方程的分量形式
\begin{equation}
F_{z\bar z}=\tfrac{\ii}{2}D_y\phi_0,\qquad
F_{yz}=\ii D_z\phi_0,\qquad
F_{y\bar z}=-\ii D_{\bar z}\phi_0,
\label{eq:bogcomponents}
\end{equation}
给出 $\partial_yA_z=-\ii D_z\phi_0$。右端正是 $A_z$ 在由 $-\ii\phi_0$ 生成的
\emph{无穷小规范变换}下的变化；因此在模掉规范变换的意义下，$E_y$ 的全纯结构
（$\bar\partial$ 算子）不随 $y$ 变。\pA\anchor{src p.128}\anchor{src p.129}
\anchor{eq 9.7}

\paragraph{第 2 步（A）：跨过 $y_0$ 时 $E_y$ 只在 $y_0$ 跳变，且挖去 $x$ 后跳变是平凡的。}
$y=y_0$ 处 $E_y$ 无定义（那里是奇异点）。但把 $E$ 限制到 $C\setminus\{x\}$ 上时，
第 1 步的论证照样适用，故限制丛的全纯类型在 $y_0$ 处也不跳。于是跳变只可能发生在
$x$ 这一点上：在 $C\setminus\{x\}$ 上 $E_+$ 与 $E_-$ 自然同构。\pA\anchor{src p.129}

第 1、2 步合起来已经给出了 Hecke 修改的\emph{全部}结构：一对丛 $E_\pm$，
一个在 $C\setminus\{x\}$ 上的同构。剩下的问题是：这个修改的\emph{类型}是什么？

\paragraph{第 3 步（A）：$G=U(1)$ 时，$L_+\cong L_-\otimes\mathcal{O}(x)^q$，且 $q=m$。}
设 $E$ 是线丛 $L$。因为 $L_+\otimes L_-^{-1}$ 在 $x$ 外平凡，它必形如
$\mathcal{O}(x)^q$（$q$ 为那个在 $x$ 外无零点无极点的截面的零点阶）。
剩下要证 $q=m$，其中 $m$ 是 't Hooft{} 算子的磁荷。原文的论证是拓扑的，我们把它写成
一次完整的同调计算：取 $y_-<y_0<y_+$、$C_\pm=\{y_\pm\}\times C$，记
$q_\pm=\oint_{C_\pm}c_1(L)$；再取包围 $p$ 的小球面 $S$，按图 14 的定向，
同调圈
\begin{equation}
D=C_+-C_--S
\label{eq:cycleD}
\end{equation}
在 $W\setminus\{p\}$ 中是边界，故
\begin{align}
0=\oint_D c_1(L)
&=\oint_{C_+}c_1(L)-\oint_{C_-}c_1(L)-\oint_S c_1(L)\notag\\
&=q_+-m-q_-,
\qquad\text{即}\qquad q_+=q_-+m .
\label{eq:homology}
\end{align}
于是 $L_+\cong L_-\otimes\mathcal{O}(x)^m$：$G=U(1)$ 的 't Hooft{} 算子 $T(m)$ 的作用
就是“用 $\mathcal{O}(x)^m$ 扭化”，与 $\mathrm{U}(1)$ Hecke 算子的标准定义一致。
\pA\anchor{src p.130}\anchor{src p.131}

这段计算值得多看一眼：它并没有用到 Bogomolny 方程的任何细节，只用了
“$'t Hooft$ 算子的奇异球面上 $\oint c_1=m$” 这一条拓扑事实，加上
“$C\setminus\{x\}$ 上同构”。换句话说，在 abelian 情形，$M3$ 是\emph{纯拓扑}的：
拓扑决定了 $T(m)$ 的作用，没有连续参数可言（$\mathcal{Y}(m)$ 是单点，与命题
\ref{prop:gl1} 一致）。非 abelian 时“远非如此”——这正是下一节要处理的。

\paragraph{第 4 步（A）：$G=U(N)$ 时，相对分解 $E=\oplus L_i$ 有
$L_i\mapsto L_i\otimes\mathcal{O}(x)^{m_i}$。}
在奇异点附近，Bogomolny 方程约化到极大环面 $T=U(1)^N$ 的方程（全纯的观点更自然：
用复化极大环面，$U(N)_{\CC}=GL(N,\CC)$）。于是可以选一个局部分解
$E|_{\text{附近}}=\oplus_{i=1}^NL_i$，使 't Hooft{} 算子在每个 $L_i$ 上的作用与 abelian 情形相同：
\begin{equation}
T(\Ldual w):\quad L_i\longmapsto L_i\otimes\mathcal{O}(x)^{m_i},
\qquad \Ldual w=(m_1,\dots,m_N).
\label{eq:UNaction}
\end{equation}
因此 $c_1(E)\mapsto c_1(E)+\sum_im_i$。\pA\anchor{src p.132}\anchor{src p.133}

\begin{origquote}
“$T(\Ldual w)$ acts by $L_i\mapsto L_i\otimes\mathcal{O}(p_0)^{m_i}$ in a local decomposition $E=\oplus L_i$---this is the basic link between 't Hooft{} and Hecke operators.”
\end{origquote}

上面这段引文是本文对 \textbf{M3} 最直接的陈述。把它与
\eqref{eq:c1shift} 一起写出来：
\begin{equation}
c_1(E)\longmapsto c_1(E)+\Bigl(\sum_i m_i\Bigr)[x],
\label{eq:c1shift}
\end{equation}
这解释了原文紧接着的一句：$E$ 的拓扑不变当且仅当 $\sum_im_i=0$，
而这又恰好是“$\Ldual G$ 的中心在 $\Ldual R$ 上平凡作用”的条件。\pA\anchor{src p.133}

\paragraph{第 5 步（A）：一般 $G$ 用权空间分解。}
对一般 $G$，取 $G$ 的表示 $R$ 与丛 $E_R=E\times_GR$；取极大环面 $T$ 与权 $w:T\to U(1)$，
按权空间分解 $E_R=\oplus_wE_{R,w}$（把 $T$ 延拓为 $x$ 附近全纯变动的极大环面团）。
由 $x$ 附近 Bogomolny 方程的 abelian 性质，相对该分解
\begin{equation}
T(\Ldual w):\quad E_{R,w}\longmapsto E_{R,w}\otimes\mathcal{O}(x)^{\langle\Ldual w,w\rangle},
\label{eq:weightaction}
\end{equation}
其中 $\langle\Ldual w,w\rangle\in\ZZ$ 是权与余权的配对。\pA\anchor{src p.134}

\paragraph{第 6 步（B）：为什么这“就是” Hecke 算子，而不只是“某个类似的动作”。}
第 4、5 步只证明了“作用的形式”与 Hecke 修改的形式相同；要断言“就是”，
还需两件事：(i) 这些修改的全体（而不是单个）组成 Hecke 算子的定义域；
(ii) 修改的类型正是 $\rho$ 的共轭类（即 GNO 荷）。原文在 §9.2、§9.3 补上 (i)：
所有可能的 $E_+$ 构成空间 $\mathcal{Y}(\Ldual w)$，且它等于仿射 Grassmann 中的一个
$G(\mathcal{O})$-轨道；在 §10.3 补上 (ii)：模空间 $\mathcal{Z}$ 与 $\mathcal{Y}$ 是同一个复流形。
我们在这里先把结论记下，证明留给 §\ref{sec:flesh-hecke-gr} 与 §\ref{sec:flesh-hecke-s10}。
\pB\anchor{src p.136；重建依据：p.146--p.148 的 Gr 分层}

\paragraph{第 7 步（B）：“类型由 GNO 荷决定”的机制。}
把 \eqref{eq:UNaction} 与 Hecke 修改的定义对照：$E_+$ 与 $E_-$ 的差被 $\mathcal{O}(x)^{m_i}$
完全刻画，而 $m_i$ 正是 $\rho(\ee^{\ii\alpha})=\mathrm{diag}(\ee^{\ii m_i\alpha})$ 的指数。
另一方面，Hecke 修改的类型按定义 \eqref{eq:doublecoset} 是双陪集，由支配余权
$(m_1,\dots,m_N)$ 刻画。两组数据逐位相等，不是“恰好同构”而是“同一个”：
$\rho$ 的 Lie 代数像就是 $z^{m_i}$ 的赋值。这一点在 §\ref{sec:flesh-hecke-gr} 写成
$g(z)=\mathrm{diag}(z^{-m_i})=\rho_{\CC}(z)$。\pB\anchor{src p.147；重建依据：p.131 的 (9.11)
与 p.147 的 (9.31) 是同一个整数 $N$-plet}

\begin{table}[htbp]
\centering\footnotesize
\caption{§9.1 论证链的七步与 provenance 分布。}
\label{tab:chain}
\begin{tabular}{@{}clll@{}}
\toprule
\tabhead{步} & \tabhead{内容} & \tabhead{级别} & \tabhead{锚点}\\
\midrule
1 & 远离 $y_0$：$E_y$ 全纯类型不随 $y$ 变 & A & p.128--129\\
2 & 只在 $y_0$ 跳变，且 $C\setminus\{x\}$ 上平凡 & A & p.129\\
3 & $U(1)$：$L_+=L_-\otimes\mathcal{O}(x)^m$ & A & p.130--131\\
4 & $U(N)$：$L_i\mapsto L_i\otimes\mathcal{O}(x)^{m_i}$ & A & p.132--133\\
5 & 一般 $G$：$E_{R,w}\mapsto E_{R,w}\otimes\mathcal{O}(x)^{\langle\Ldual w,w\rangle}$ & A & p.134\\
6 & 修改的全体 $=\mathcal{Y}(\Ldual w)$ $=$ Gr 轨道 & B & p.136, p.146\\
7 & 类型 $=$ GNO 荷：$\rho$ 的指数就是 $m_i$ & B & p.147\\
\bottomrule
\end{tabular}
\end{table}

\subsubsection{M3 的精确陈述}

把七步合起来，得到本节的主定理。注意它是\emph{物理陈述}：它的“证明”是把
两个模空间认出来（§10.3），而不是形式推导。

\begin{definition}[M3 的精确版本]\label{def:m3}
在 $M=\RR\times I\times C$、$\Psi=0$、$W=I\times C$ 的设定下，于 $p=y_0\times x$
处插入型 $\Ldual w$ 的静态 't Hooft{} 算子。设 $y_-<y_0<y_+$，$E_\pm=E|_{\{y_\pm\}\times C}$。
则 $E_+$ 与 $E_-$ 在 $C\setminus\{x\}$ 上自然同构，且这个同构在 $x$ 处恰好实现型
$\Ldual w$ 的 Hecke 修改：
\begin{equation}
E_+\cong E_-\ \text{在}\ C\setminus\{x\}\ \text{上同构},\qquad
E_+\big|_{\text{附近}}=\bigoplus_i L_i\otimes\mathcal{O}(x)^{m_i},
\label{eq:m3}
\end{equation}
其中 $L_i$ 是 $E_-$ 在 $x$ 附近的局部分解。并且这个修改的\emph{型}
（双陪集 \eqref{eq:doublecoset}）等于 't Hooft{} 算子定义中同态 $\rho$ 的共轭类。
\end{definition}

作为“算子”，Hecke 算子的完整定义是把 \eqref{eq:m3} 的所有可能输出求和（或积分）：
对 $\mathrm{Bun}_G$ 上的态（函数、层、或 A-模型态空间中的元素）$\Psi$，
\begin{equation}
(\mathbb{H}_{\Ldual w}\Psi)(E)=\int_{\mathcal{Y}(\Ldual w;x)}\Psi(E')\,D(\text{参数}),
\label{eq:heckeop}
\end{equation}
其中积分跑遍 $E$ 的全部型 $\Ldual w$ 修改（在紧化空间上积分；见 §\ref{sec:flesh-hecke-gr}）。
物理上，\eqref{eq:heckeop} 就是“在 $x$ 处插一个 't Hooft{} 算子，
然后对 $x$ 附近所有允许的场构型求和”。\pB\anchor{src p.149；重建依据：p.171 的
“态空间 $=$ $\mathcal{Z}$ 的上同调”与 p.162 的“$\mathcal{Z}'=$ 修改空间”}

\subsubsection{同一个 $\rho$ 的三次出现：本节的枢纽}

\begin{keyidea}
\textbf{三重巧合（其实不是巧合）。}同一个同态 $\rho:U(1)\to G$ 在本节出现三次：
作为\emph{奇异性}（它规定沿奇异线的曲率行为，\eqref{eq:Fdiag}）、
作为\emph{GNO 荷}（其共轭类由 $\Ldual G$ 的最高权分类，即 't Hooft{} 算子的标签）、
作为\emph{仿射 Grassmann 的点}（$\rho_{\CC}:\CC^*\to G_{\CC}$，$z\mapsto\mathrm{diag}(z^{m_i})$）。
M3 的内容就是：这三者是同一个东西的三种写法。\pA\anchor{src p.132}\anchor{src p.84}\anchor{src p.147}
\end{keyidea}

\subsubsection{排除错误直觉 4：“算子”二字不能按字面读}

原文在 §9 开头有一句极重要的话，我们在前文引过它的意思：这里的线“算子”
不是通常意义下“作用于一个既有向量空间”的算子，它必须被纳入\emph{物理超对称
态空间的定义}之中。\pA\anchor{src p.123}

具体地说：在二维 A-模型里，插入 't Hooft{} 算子改变的是\emph{边界条件}（膜的选择），
从而改变了我们量子化的对象——态空间是 Bogomolny 模空间 $\mathcal{Z}$ 的上同调，
而不是某个固定 Hilbert 空间的子空间。\pA\anchor{src p.171} 把 't Hooft{} 算子想成
“作用在态上的矩阵”，会在两个地方立刻失败：它没有“与态的内积”可言，
也不能脱离膜的选择单独定义。正确的图像是“一个对应（correspondence）”：
$\mathcal{Z}$ 是 $E_-\times E_+$ 上的子簇，算子是把一侧的类推到另一侧。这正是
\eqref{eq:heckeop} 的几何版本。\pB\anchor{src p.150；重建依据：p.151 的 Hecke 对应 $Q$}

\begin{kcard}{排除错误直觉 5：三个对偶，别混}
\anchor{src p.84；src p.123；src p.147}
\textbf{错误理解：}“S-对偶把 Wilson 变成 't Hooft，而 't Hooft 又是 Hecke，
所以 Hecke 算子就是 Wilson 算子的 S-对偶。”

\textbf{正确图像：}本文涉及三个\emph{互不相同}的对偶，下表给出它们各自交换什么、
在本文哪里被使用、以及\textbf{不能}从中推出什么。
\begin{center}\footnotesize
\begin{tabular}{@{}p{2.0cm}p{4.5cm}p{4.3cm}p{3.4cm}@{}}
\toprule
\tabhead{对偶} & \tabhead{交换什么} & \tabhead{在本文的用途} & \tabhead{不能推出}\\
\midrule
S-对偶 & $\tau\leftrightarrow-1/\tau$；Wilson $\leftrightarrow$ 't Hooft；
$\Psi=\infty\leftrightarrow\Psi=0$ & 说明 $\Psi=0$ 处\emph{应当}存在与 Wilson 对偶的算子（§6.2） & Hecke 算子的类型或修改空间\\
GNO 对偶 & $G\leftrightarrow\Ldual G$；表示 $\leftrightarrow$ 余权 &
分类 't Hooft{} 算子：$T(\Ldual R)$ 由 $\Ldual G$ 的表示标号（§6.2） & 算子的几何作用\\
几何 Langlands & $\MH(G)$ 的 A-膜 $\leftrightarrow\MH(\Ldual G)$ 的 B-膜 &
Hecke 算子作用在膜上 $=$ 丛的 Hecke 修改（§9--§10） & 算子乘积的具体系数\\
\bottomrule
\end{tabular}
\end{center}
M3 属于第三行。它把\emph{第一行给出的算子}与\emph{几何 Langlands 的} Hecke 算子
接起来；接法用的是第二行的标签。三者缺一不可，但任何一者单独都不等于 M3。
\end{kcard}

\paragraph{自检问题。}
(1) 在 \eqref{eq:homology} 中，如果 $S$ 的定向反过来，$q=m$ 会变成什么？
这与 \eqref{eq:flux1} 的归一化有什么关系？
(2) 为什么第 1 步的“全纯类型不变”需要“模掉规范变换”才成立？
如果不模，$E_y$ 到底是什么？
(3) 用 \eqref{eq:weightaction} 计算：若 $\Ldual w$ 是权 $\Ldual R$ 的最高权，
$E_{R,w}$ 在 $T(\Ldual w)$ 下被张量多少次 $\mathcal{O}(x)$？给出一个非零例子。


\subsection{第四层：仿射 Grassmann 与 Schubert 分层（§9.3，任务第 5 点）}
\label{sec:flesh-hecke-gr}

\subsubsection{定义：修改空间是无穷维流形的有限维子簇}

上一节把 Hecke 修改的类型说成双陪集 \eqref{eq:doublecoset}，但那只给出了\emph{类型}，
没有给出\emph{空间}。要理解修改空间的几何（从而理解 't Hooft{} 算子路径积分积什么），
需要一个把“所有修改”放在一起的框架。这个框架就是
\term{affine Grassmannian}{仿射 Grassmann 流形}。

取 $C=\PP^1$、$x=0$，局部坐标 $z$。设 $E=\mathcal{O}^{\oplus N}$ 平凡，并固定
它在穿孔圆盘上的平凡化。所有可能的 Hecke 修改（即所有二元组 $(E',\beta)$，模去延拓过
$0$ 的同构）构成的空间是
\begin{equation}
\mathrm{Gr}_N=GL(N,\CC((z)))/GL(N,\CC[[z]]).
\label{eq:grN}
\end{equation}
一般的 $G$ 版本是
\begin{equation}
\mathrm{Gr}_G=G_{\CC}((z))/G_{\CC}[[z]].
\label{eq:grG}
\end{equation}
\pA\anchor{src p.146}

\begin{remark}[为什么这个商描述的对子是 $(E',\beta)$]
$\mathrm{Gr}_N$ 的元素是 $\CC((z))$ 中可逆矩阵 $g$ 的陪集。给定 $g$，用 $g$ 把平凡丛
在穿孔圆盘上的两个平凡化粘起来，就得到曲线上一个新的丛 $E'$，而“粘合矩阵恰好是 $g$”
就是 $E'$ 与平凡丛在 $0$ 外的同构 $\beta$。右乘 $GL(N,\CC[[z]])$ 相当于在圆盘上换
$E'$ 的平凡化（不改变 $E'$），左乘相当于换 $E$ 的平凡化。所以 $\mathrm{Gr}_N$
的点是“$E'$ 加上 $\beta$”，只模掉两边的平凡化选择。\pA\anchor{src p.146}
\end{remark}

原文 §9.3 给了一个更几何的版本（我们视为 $B$ 级重建，因为它依赖引文 [122] 的定理）：
设 $X$ 是 $C^*$ 上全纯、并可亚纯延拓到 $\PP^1$ 的函数环，$\mathcal{O}$ 是 $\CC[[z]]$
的对应物。则
\begin{equation}
\bigl\{(E,\sigma)\bigr\}/\sim\ \cong\ GL(N,X)/GL(N,\mathcal{O})\ \cong\ \mathrm{Gr}_N,
\label{eq:grX}
\end{equation}
其中 $E$ 是 $\PP^1$ 上秩 $N$ 全纯丛，$\sigma$ 是 $E$ 在 $U_\infty=\PP^1\setminus\{0\}$
上的平凡化。论证是：取 $E$ 在 $U_0=\PP^1\setminus\{\infty\}\cong\CC$ 上的平凡化 $\sigma'$，
则 $\sigma$ 与 $\sigma'$ 在重叠区域 $U_0\cap U_\infty=\CC^*$ 上的转移函数是
$g(z)\in GL(N,X)$；换 $\sigma'$ 相当于右乘 $GL(N,\mathcal{O})$。于是
\eqref{eq:grX} 成立。\pB\anchor{src p.146}

关键的一步是：\textbf{这些无穷维的东西，被修改的“型”切成有限维的块}。
给定支配余权 $\Ldual w=(m_1,\dots,m_N)$（排序 $m_1\ge\cdots\ge m_N$），
对应的 $GL(N,\mathcal{O})$-轨道是
\begin{equation}
\mathcal{Y}(\Ldual w;0)=GL(N,\mathcal{O})\cdot
\mathrm{diag}(z^{-m_1},\dots,z^{-m_N})
\ \subset\ \mathrm{Gr}_N,
\label{eq:orbit}
\end{equation}
它就是型 $\Ldual w$ 的 Hecke 修改空间。\pA\anchor{src p.147}\anchor{eq 9.31}

\subsubsection{枢纽：$\mathrm{diag}(z^{-m_i})=\rho_{\CC}(z)$}

现在写下令 M3 从“巧合”变成“恒等式”的那一行。把 \eqref{eq:orbit} 的矩阵
与 't Hooft{} 算子的定义 \eqref{eq:rhoUN} 并排：
\begin{equation}
\underbrace{\rho(\ee^{\ii\alpha})=\mathrm{diag}(\ee^{\ii m_1\alpha},\dots,\ee^{\ii m_N\alpha})}
_{\text{'t Hooft 算子的奇异性 / GNO 荷}}
\qquad\xrightarrow{\ \ee^{\ii\alpha}\leftrightarrow z\ } \qquad
\underbrace{g(z)=\mathrm{diag}(z^{-m_1},\dots,z^{-m_N})}
_{\text{仿射 Grassmann 的点 / Hecke 修改的类型}}.
\label{eq:rhoC}
\end{equation}
左边是物理定义的算子标签，右边是数学定义的修改类型；中间的箭头只是把 $U(1)$ 的
参数 $\ee^{\ii\alpha}$ 解析延拓到 $\CC^*$ 的坐标 $z$。\pA\anchor{src p.147}
\anchor{eq 9.32}

原文对这一步的说法是：\eqref{eq:orbit} 中的 $g(z)$ 是 $\CC^*\to GL(N,\CC)$ 的同态，
它“实为” $\rho:U(1)\to U(N)$ 的复化，而 $\rho$ 正是 't Hooft{} 算子
$T(m_1,\dots,m_N)$ 定义中的同态。\pA\anchor{src p.147}

\begin{keyidea}
\textbf{M3 的组合核心。}'t Hooft{} 算子的 GNO 标签 $\Ldual w$ 是 $\Ldual G$ 的最高权；
Hecke 修改的类型也是 $\Ldual G$ 的最高权。这不是“两个同构的对象恰好同构”，
而是“同一组整数 $m_1,\dots,m_N$ 在两处各写了一遍”：
一处写成 $\rho$ 在 $\mathfrak{u}(1)$ 上的像（李代数），一处写成 $\rho_{\CC}$ 在
$\CC((z))$ 中的像（群）。原文用同一个符号 $\Ldual w$ 记这两样东西，正是这个意思。\pA\anchor{src p.147}
\end{keyidea}

对带中心的一般群，这个对应还要加一个“扭”的成分。以 $G=P\,SU(N)$ 为例
（$\Ldual G=SU(N)$）：$P\,SU(N)$ 丛没有对应的秩 $N$ 向量丛，$m_i$ 允许取分数，
但要求 $m_i-m_j\in\ZZ$。修改空间落在“扭”的仿射 Grassmann 版本里，
$z^{-m_i}$ 的分数指数彼此相消，作用到任何 $P\,SU(N)$ 的表示丛上都没有分数指数。
原文用 $SO(3)$ 的伴随丛 $L^2\oplus\mathcal{O}\oplus L^{-2}$ 演示了这一点。\pA\anchor{src p.135}

\subsubsection{两种分层，以及“紧化”到底紧化什么}

$\mathrm{Gr}_N$ 上有两套完全不同的分层，初学时极易混淆，而 M3 的“修改空间”用的是
其中一套。

\paragraph{分层 A（G(O)-轨道分层，有限维）：}
按 $GL(N,\mathcal{O})$-轨道分层，每个轨道是一个 $\mathcal{Y}(\Ldual w)$，
由整数 $N$-plet 模置换标号。\textbf{修改空间用的是这一套。}
每一层\emph{不是}紧的（除少数例外），其闭包是 Schubert 簇。
$\mathrm{Gr}_N$ 中每个 $GL(N,\mathcal{O})$-轨道都形如 \eqref{eq:orbit}
（原文引 [122] 的定理）。\pA\anchor{src p.147}

\paragraph{分层 B（Grothendieck 分层，无穷维）：}
由 Grothendieck 定理，$\PP^1$ 上秩 $N$ 全纯丛
\begin{equation}
E\cong\bigoplus_{i=1}^N\mathcal{O}(x)^{k_i}\qquad(k_1\ge\cdots\ge k_N).
\label{eq:grothendieck}
\end{equation}
按“$E$ 的同构类”分层（即固定 $E$ 而让平凡化 $\sigma$ 变），每层显然无穷维：
固定丛的平凡化有无穷多个。\pA\anchor{src p.147}

\begin{kcard}{排除错误直觉 6：$\mathrm{Gr}$ 不是“丛的模空间”}
\anchor{src p.146；src p.147}
\textbf{错误理解：}仿射 Grassmann 流形是“$\PP^1$ 上 $G$-丛的模空间”，
所以修改空间是“丛的空间”的一个子簇。

\textbf{为什么错：}若真是丛的模空间，它就该按同构类分层（分层 B），而分层 B 是无穷维的。
$\mathrm{Gr}_N$ 的点其实是“丛 $+$ 在 $x$ 外的平凡化”，即
$(E,\sigma)$，其维数无穷但轨道有限维。修改空间 $\mathcal{Y}(\Ldual w)$ 只是其中一个
\emph{有限维} $GL(N,\mathcal{O})$-轨道，不是丛模空间的子簇。\pA\anchor{src p.146}

\textbf{正确图像：}$\mathrm{Gr}_N$ 是“局部修改的数据”的空间。它的有限维子簇
（Schubert 簇）才是“修改空间”。丛本身的模空间 $\mathrm{Bun}_G$ 是另一个
（无穷维的）对象，Hecke 算子是 $\mathrm{Bun}_G$ 上的对应，而 't Hooft{} 算子的路径积分
在\emph{纤维} $\mathcal{Y}$ 上做。
\end{kcard}

\paragraph{紧化。}按定义 \eqref{eq:orbit}，$\mathcal{Y}(\Ldual w)$ 是一族光滑点的集合，
它一般\emph{不是}闭的。自然的紧化是取它在某个有限维截断中的闭包：
\begin{equation}
\overline{\mathcal{Y}}(\Ldual w)=\overline{GL(N,\mathcal{O})\cdot g_{\Ldual w}}
\quad\subset\quad \mathrm{Gr}_N(i),\qquad
\mathrm{Gr}_N(i)=\{g:\ \text{极点阶}\le i\}.
\label{eq:closure}
\end{equation}
闭包一般是奇异簇，它包含 $\mathcal{Y}(\Ldual w)$ 本身，以及若干\emph{低维}轨道。\pA\anchor{src p.147}
哪些低维轨道会出现？原文的回答是：所有\term{associated weight}{相伴权}型的修改。
称 $\Ldual w'$ 是 $\Ldual w$ 的 associated 权，若 $\Ldual w'$ 是 $\Ldual w$ 的一个权
但非最高权，等价地 $\Ldual w-\Ldual w'$ 是伴随表示的\emph{非零}支配权。
对 $U(N)$，这等价于：存在 $i,j$ 使 $m_i\ge m_j+2$，把 $(m_i,m_j)$ 换成
$(m_i-1,m_j+1)$。\pA\anchor{src p.143}

\begin{figure}[htbp]
\centering
\begin{tikzpicture}[scale=1.0,>=Stealth,font=\small]
 % ---- left: the truncated affine Grassmannian ----
 \draw[thick,kwblue,fill=kwblue!3] (-1.2,0) ellipse (3.3 and 2.05);
 \draw[thick,kwred,dashed,fill=kwred!4] (-1.5,0.45) ellipse (2.55 and 1.15);
 \draw[thick,kwred] (-2.2,0.70) ellipse (0.95 and 0.40);
 \draw[thick,kwgreen,fill=kwgreen!12] (-1.9,-0.60) ellipse (1.05 and 0.42);
 \draw[thick,kwred,fill=kwred!12] (0.25,-0.50) ellipse (0.55 and 0.24);
 \fill (1.15,-1.45) circle (2pt);
 \draw[->,kwgray,dashed] (-1.20,0.32) -- (-0.20,-0.30);
 \node[font=\scriptsize,kwgray] at (-1.55,-0.02) {$(2,0)\to(1,1)$};
 % ---- right: legend ----
 \begin{scope}[shift={(2.6,1.75)}]
   \draw[thick,kwblue,fill=kwblue!3] (0,0) ellipse (0.22 and 0.12);
   \node[font=\footnotesize,anchor=west] at (0.38,0)
     {$\mathrm{Gr}_N(i)$：$z=0$ 处极点阶 $\le i$ 的截断};
   \draw[thick,kwred,dashed,fill=kwred!4] (0,-0.62) ellipse (0.22 and 0.12);
   \node[font=\footnotesize,anchor=west] at (0.38,-0.62)
     {边界层：associated 权型的修改};
   \draw[thick,kwred] (0,-1.24) ellipse (0.22 and 0.12);
   \node[font=\footnotesize,anchor=west] at (0.38,-1.24)
     {$\mathcal{Y}((2,0))$：奇异的 Schubert 簇};
   \draw[thick,kwgreen,fill=kwgreen!12] (0,-1.86) ellipse (0.22 and 0.12);
   \node[font=\footnotesize,anchor=west] at (0.38,-1.86)
     {$\mathcal{Y}((1,0))\cong\PP^1$：极小、闭、光滑};
   \draw[thick,kwred,fill=kwred!12] (0,-2.48) ellipse (0.22 and 0.12);
   \node[font=\footnotesize,anchor=west] at (0.38,-2.48)
     {$\mathcal{Y}((1,1))$：边界层上的退化};
   \fill (0,-3.10) circle (2pt);
   \node[font=\footnotesize,anchor=west] at (0.38,-3.10)
     {$\mathcal{Y}(0)$：平凡丛（一点）};
 \end{scope}
\end{tikzpicture}
\caption{$GL_2$ 的仿射 Grassmann 中被修改空间占据的部分：$\mathcal{Y}((1,0))$ 是极小表示型，
闭且光滑；$\overline{\mathcal{Y}}((2,0))$ 的边界含 $\mathcal{Y}((1,1))$（以及平凡型）。
一般地，闭包的边界层恰是 associated 权型的修改空间。}
\label{fig:schubert}
\end{figure}

\subsubsection{维数公式：原文省略的推导}

修改空间的维数是本节唯一需要“算”的量；它后来在 §10.2 被指标理论独立验证
（\eqref{eq:pauly}）。原文把推导压缩成三行，我们补全。

\begin{proposition}[修改空间的维数]\label{prop:dim}
对 $G=U(N)$（或 $SU(N)$、$P\,SU(N)$），型 $\Ldual w=(a_1,\dots,a_N)$（$a_1\ge\cdots\ge a_N$）
的修改空间复维为
\begin{equation}
\Delta_{\Ldual w}=\dim_{\CC}\mathcal{Y}(\Ldual w)=\sum_{i<j}(a_i-a_j).
\label{eq:dimformula}
\end{equation}
\end{proposition}

\begin{proof}[推导]
取 $N$ 个截面 $s_i$ 在 $0$ 附近生成平凡丛 $E_-=\mathcal{O}^{\oplus N}$
（即 $s_i(0)$ 线性无关）。型 $\Ldual w$ 的修改空间中的一般点可写成
\begin{equation}
q=\sum_{i=1}^N z^{-a_i}g_i(z)\,s_i(z),
\label{eq:generalpoint}
\end{equation}
其中 $g_i$ 在 $0$ 全纯。这不是唯一的写法：作变换
\begin{equation}
s_i\longmapsto s_i+\sum_{j\le i}h_js_j+\sum_{j>i}z^{a_i-a_j}h_js_j
\label{eq:gaugetransf}
\end{equation}
（$h_j$ 全纯）不改变 $E_+$，且可被吸收进 \eqref{eq:generalpoint} 中 $g_i$ 的重定义。
把 $s_i$ 作 Taylor 展开，“不能被 \eqref{eq:gaugetransf} 消去”的系数，
恰是那些在 $s_i$ 的展开里出现、阶数低于 $z^{a_i-a_j}$ 的项；对固定的 $i$，
这样的系数有 $\sum_{j>i}(a_i-a_j)$ 个。对 $i$ 求和得
\begin{equation}
\Delta_{\Ldual w}
=\sum_{i=1}^N\sum_{j=i+1}^N(a_i-a_j)
=\sum_{i<j}(a_i-a_j).
\label{eq:dimderive}
\end{equation}
\end{proof}

\begin{example}[两个检验]
对 $\Ldual w=(1,0)$：$\Delta=1$，与 $\mathcal{Y}((1,0))=\PP^1$ 的复维 $1$ 一致。
对 $\Ldual w=(a,0)$：$\Delta=a$，与 §9.2 的“型 $(a,0)$ 的修改空间复维 $a$” 一致
（那个例子是逐项数极点系数得到的，见 \eqref{eq:y20} 的一般化）。\pA\anchor{src p.143}
\end{example}

\begin{remark}[为什么维数公式必须与指标理论一致]
\eqref{eq:dimformula} 是纯代数几何的计数；\eqref{eq:pauly} 是四维等变瞬子的指标计算。
两者相等这件事，是 M3 的\emph{独立检验}：几何上数出来的参数个数，恰好等于物理上
Bogomolny 模空间的（虚拟）维数。如果它们不等，M3 就是错的。\pB\anchor{src p.155}
\end{remark}

\paragraph{自检问题。}
(1) 在 \eqref{eq:grX} 中，如果把 $\sigma$ 换成 $E$ 在 $U_0$ 上的平凡化而不是
$U_\infty$ 上的，得到的空间还是 $\mathrm{Gr}_N$ 吗？差在哪里？
(2) 为什么分层 A 是有限维的，而分层 B 是无穷维的？用“固定什么、让什么变”回答。
(3) 对 $G=SU(2)$，型 $(1,-1)$（伴随表示）的修改空间是什么？它的维数与
\eqref{eq:dimformula} 是否一致？


\subsection{第五层：Bogomolny 方程给出的精确框架（§10.1--§10.2）}
\label{sec:flesh-hecke-s10}

前三层已经把话说完了：'t Hooft{} 算子的奇异性由 GNO 权标号，跨越奇异性时丛做 Hecke 修改，
两者用同一组整数标号。但“说完”不等于“证完”。前三层的论证里至少有四个地方欠账：
(a) 为什么两个模空间\emph{相等}而不只是\emph{维数相同}？
(b) 修改空间一般非紧，而路径积分遍历\emph{所有}场构型，非紧部分贡献什么？
(c) 修改空间一般奇异，奇异点上的上同调（物理态空间）怎么定义？
(d) 两个 't Hooft{} 算子靠近时发生什么（OPE）？\pC

§10 的全部技术内容就是在回答这四个问题。本层按原文顺序走：§10.1 边界条件与指标，
§10.2 单极子冒泡，§10.3 Kähler 商，§10.4 OPE 与 Hecke 代数，§10.5 扩展 Bogomolny 方程。

\subsubsection{§10.1：把“算子作用”翻译成“边值问题”}

\paragraph{动机。}取 $W=I\times C$（$I$ 为区间），在内部若干点 $p_i$ 插入 't Hooft{} 算子。
我们想研究：给定左端 $E_-$，右端能出现什么 $E_+$。要把这变成一个\emph{椭圆}边值问题，
必须对 $(A,\varphi_0)$ 在两端各给一半边界条件。

\paragraph{推导：为什么是“一端 Dirichlet、一端 Neumann”。}
Bogomolny 方程 \eqref{eq:bogomolny} 是二阶方程，未知量是 $(A,\varphi_0)$。若两端都把 $A$
定死，方程会过度约束（一般无解）；若两端都放开，解不唯一、模空间无穷维，指标无意义。
最基本的椭圆边界条件是：
\begin{itemize}
\item 在一端给定 $A$ 的边界值（\term{Dirichlet}{狄利克雷}），让 $\varphi_0$ 的边界值自由；
  方程随即决定 $\varphi_0$ 的法向导数。若该端 $A_-$ 平坦，则要求 $\partial_n\varphi_0=0$。
\item 在另一端设 $\varphi_0=0$，让 $A$ 自由；方程随即要求 $A$ 满足协变 Neumann 条件
  $\iota_nF=0$（曲率与法向量缩并为零）。
\end{itemize}
这两条正是原文 Fig.~15(a) 的边界条件：左端“$A$ 定、$\varphi_0$ 不定”，
右端“$\varphi_0=0$、$A$ 不定”。\pA\anchor{src p.149}\anchor{src p.150}

\begin{figure}[htbp]
\centering
\begin{tikzpicture}[scale=1.0,>=Stealth,font=\small]
 % ---- (a) W = I x C: holomorphic type jumps once ----
 \node[font=\footnotesize] at (-2.6,3.0) {给定 $A_-$，$\varphi_0$ 自由};
 \node[font=\footnotesize] at (2.6,3.0) {$\varphi_0=0$，$\iota_nF=0$};
 \draw[thick,fill=kwblue!8] (-3.7,2.15) rectangle (-1.5,2.6);
 \node[font=\footnotesize] at (-2.6,2.37) {$C_-$};
 \draw[thick,fill=kwgreen!10] (1.5,2.15) rectangle (3.7,2.6);
 \node[font=\footnotesize] at (2.6,2.37) {$C_+$};
 \draw[->,thick] (-2.6,2.1) -- (-2.6,1.25) node[midway,left,font=\footnotesize] {$I$};
 \draw[->,thick] (2.6,2.1) -- (2.6,1.25);
 \draw[very thick,kwblue] (-3.5,1.0) -- (-0.12,1.0);
 \draw[very thick,kwred] (0.12,1.0) -- (3.5,1.0);
 \fill (0,1.0) circle (2.6pt);
 \node[font=\footnotesize,anchor=west] at (0.28,0.6) {$p=y_0\times x$};
 \node[font=\footnotesize,kwblue] at (-3.5,0.6) {$E_-$};
 \node[font=\footnotesize,kwred] at (3.5,0.6) {$E_+$};
 \node[font=\footnotesize,align=center,text width=8.6cm] at (0,0.0)
   {跨过 $p$（'t Hooft{} 算子位置）时 $E|_{\{y\}\times C}$ 的全纯型跳变一次};
 % ---- (b) local: isomorphic away from x ----
 \begin{scope}[shift={(5.4,1.85)}]
   \draw[thick] (0,0) circle (0.8);
   \node[font=\footnotesize] at (0,-1.15) {$C$（或 $\PP^1$）};
   \fill (0,0) circle (2.2pt);
   \node[font=\footnotesize,above right] at (0.06,0.04) {$x$};
   \draw[thick,dashed] (0,0) circle (0.3);
   \node[font=\footnotesize,anchor=west] at (1.0,0.25) {$E_+\cong E_-\ \text{on}\ C\setminus\{x\}$};
   \node[font=\footnotesize,anchor=west] at (1.0,-0.2) {差别只在 $x$ 附近一个局部数据};
 \end{scope}
\end{tikzpicture}
\caption{局部修改的两个视图。左：$W=I\times C$ 中，$E$ 的全纯型沿 $y$ 传输，在跨过
't Hooft{} 算子位置 $p$ 时跳变一次。右：“局部”的准确含义——两条丛在挖去 $x$ 后
同构，全部差别集中在 $x$ 的任一邻域里。}
\label{fig:local}
\end{figure}

\paragraph{指标理论：能先算出维数。}线性化 \eqref{eq:bogomolny} 得椭圆复形
\begin{equation}
0\to\Omega^0(\mathrm{ad}\,E)\xrightarrow{d_1}
\Omega^1(\mathrm{ad}\,E)\oplus\Omega^0(\mathrm{ad}\,E)\xrightarrow{d_2}
\Omega^2(\mathrm{ad}\,E)\to0,
\label{eq:virtual}
\end{equation}
模空间 $Z$ 的切空间是 $H^1=\ker d_2/\operatorname{im}d_1$，其虚拟复维等于复形指标。
在原文的边界条件下 $H^0=0$（规范变换在一端平凡，故无协变常截面），
$H^2=0$（虚拟维非负，见下），故光滑点处维数就是指标。\pA\anchor{src p.152}

\paragraph{无算子时维数为零。}这是后文一切比较的基准，且原文给了完整论证：若
$F-\star D\varphi_0=0$，则
\begin{equation}
0=\int_W\Tr(F-\star D\varphi_0)\wedge\star(F-\star D\varphi_0)
=-\int_W\Tr\bigl(F\wedge\star F+D\varphi_0\wedge\star D\varphi_0\bigr)
=-2\int_{\partial W}\Tr\varphi_0F,
\label{eq:vanishing}
\end{equation}
最后一步分部积分。边界项在两端都消失：一端 $\varphi_0=0$，另一端 $F=0$。
左端半正定，故 $F=D\varphi_0=0$：解必从 $C$ 拉回，$E_+\cong E_-$，$Z$ 是一点。
\pA\anchor{src p.153}

\paragraph{Pauly 的移动量。}每个 't Hooft{} 算子 $T(\Ldual R)$ 的奇异性使虚拟维移动
$\Delta_{\Ldual R}$，移动量只依赖表示 $\Ldual R$。对 $G=P\,SU(2)$、$\Ldual G=SU(2)$，
最高权 $\Ldual w=(a/2,-a/2)$ 的算子使虚拟维移动 $a$。一般群的公式是
\begin{equation}
\Delta_\rho=\sum_k c_k,
\qquad
\rho:\ U(1)\to G\ \text{在}\ \mathrm{ad}(\widetilde E)_0\ \text{上的权分解}\
=\sum_k\bigl(\ee^{\ii c_k\theta}+\ee^{-\ii c_k\theta}\bigr),\ c_k>0,
\label{eq:pauly}
\end{equation}
即贡献等于固定点纤维上伴随作用的\emph{正权之和}。这来自 $F=U(1)$-等变瞬子的等变指标定理：
固定点对指标的贡献是 $\mathrm{ad}(\widetilde E)_0$ 上取迹，而 $\mathrm{ad}$ 表示是实的，
故非零权成对出现。\pA\anchor{src p.154}\anchor{src p.155}

多个算子时贡献相加：
\begin{equation}
\dim Z(\rho_1,\dots,\rho_n)=\sum_{i=1}^n\Delta_{\rho_i},
\label{eq:dimsum}
\end{equation}
因为指标是局部的（切除性质）。\pA\anchor{src p.155}

\begin{remark}[与 \eqref{eq:dimformula} 的对照]
\eqref{eq:pauly} 是\emph{四维}等变瞬子的指标计算，\eqref{eq:dimformula} 是\emph{一维}
（$C$ 上）的极点计数。两者相等是 M3 的一个独立检验：数参数的两种完全不同的方式
给出同一个数。原文对 $U(N)$ 的情形明确对照了这两式。\pA\anchor{src p.155}
\end{remark}

\subsubsection{§10.2：单极子冒泡——修改空间的边界从哪来}

\paragraph{问题。}$\mathcal{Y}(\Ldual w)$ 一般非紧（§\ref{sec:flesh-hecke-gr} 的
\eqref{eq:closure}）。它的紧化要把哪些点加进来？§9.2 的回答是“associated 权型的修改”，
但那只是从代数几何看出来的。物理上这些边界点是什么？

\paragraph{答案：一单位磁荷缩成一点并消失。}
把 $W=I\times C$ 局部看成 $R^3$，'t Hooft{} 算子附近的奇异 Bogomolny 解对应
$C^2\cong R^4$ 上对 $F=U(1)$ 作用
\begin{equation}
(z_1,z_2)\longmapsto(\ee^{\ii\theta}z_1,\ee^{\ii\theta}z_2)
\label{eq:faction}
\end{equation}
不变的瞬子（Kronheimer）。$F$ 作用在 $C^2$ 的原点有固定点，对应 $R^3$ 的原点，
也就是 't Hooft{} 算子的位置；$F$ 在固定点纤维 $\widetilde E_0$ 上的作用由同态
$\rho:F\to G$ 给出——这正是 't Hooft{} 算子的 GNO 权。\pA\anchor{src p.154}

\paragraph{ADHM 计算。}$R^4$ 上瞬子方程尺度不变，所以瞬子可以“缩成一点”：
这就是冒泡（bubbling）。等变情形下缩的那点自动是 $C^2$ 的原点。用 ADHM 复形
\begin{equation}
0\to H\xrightarrow{\alpha(z)}H\otimes U\oplus E_\infty
\xrightarrow{\beta(z)}H\to0,\qquad \beta(z)\alpha(z)=0,
\label{eq:adhm}
\end{equation}
（$H$ 为 $k$ 维空间、$E_\infty$ 为无穷远纤维、$U$ 为旋量空间）描述瞬子，
则在原点的纤维 $E_0$ 是 $z=0$ 处该复形的第一上同调。对只有 $H^1$ 非零的等变复形，
取特征标得
\begin{equation}
\chi(E_0)=\chi(E_\infty)+\chi(H)\bigl(\ee^{\ii\theta}+\ee^{-\ii\theta}-2\bigr).
\label{eq:chiE0}
\end{equation}
\pA\anchor{src p.156}\anchor{src p.157}

\paragraph{冒泡的代数。}冒泡就是让复形 \eqref{eq:adhm} 直和分解出一个\emph{零复形}
$0\to H_1\to H_1\otimes U\to H_1\to0$（它对上同调无贡献，但把瞬子数降低了
$\dim H_1$）。比较冒泡前后的特征标：
\begin{equation}
\chi(E'_0)=\chi(E_0)+\chi(H_1)\bigl(2-\ee^{\ii\theta}-\ee^{-\ii\theta}\bigr).
\label{eq:bubble}
\end{equation}
若 $H_1$ 是一维、特征标为 $\ee^{\ii m\theta}$，则 $E_0$ 的权表必须含 $m+1$ 与 $m-1$，
即
\begin{equation}
(m+1,\ m-1)\ \rightsquigarrow\ (m,\ m).
\label{eq:weights}
\end{equation}
一般地，$H_1$ 的权为 $m,m+1,\dots,m+s-1$ 时，移动是
$(m+s,\,m-1)\rightsquigarrow(m+s-1,\,m)$。所有 $SU(N)$ 的预期退化都是这类移动的复合；
一般 $G$ 由嵌入各 $SU(2)$ 子群的冒泡生成。\pA\anchor{src p.157}\anchor{src p.158}

\begin{figure}[htbp]
\centering
\begin{tikzpicture}[scale=1.0,>=Stealth,font=\small]
 % (a) physical picture
 \begin{scope}[shift={(-4.9,0)}]
   \draw[thick,fill=kwblue!4] (0,0) rectangle (4.0,2.5);
   \node[font=\footnotesize,anchor=north west] at (0.08,2.42) {(a) 物理图像};
   \fill (2.0,1.55) circle (2.2pt);
   \node[font=\footnotesize,above right] at (2.06,1.6) {$p$};
   \draw[thick,kwred,dashed] (2.0,1.55) circle (0.45);
   \fill[kwred!30] (2.0,1.55) circle (0.26);
   \draw[->,kwred,thick] (2.62,2.05) -- (2.22,1.72);
   \node[font=\footnotesize,kwred,align=center,text width=3.4cm] at (2.0,0.62)
     {一单位 GNO 荷\\收缩到 $p$ 并消失};
 \end{scope}
 % (b) boundary strata
 \begin{scope}[shift={(0.6,0)}]
   \node[font=\footnotesize,anchor=north west] at (0,2.6) {(b) 模空间边界层};
   \draw[thick,kwblue,fill=kwblue!3] (2.2,1.05) ellipse (1.9 and 0.95);
   \draw[thick,kwgreen,fill=kwgreen!12] (2.2,0.50) ellipse (1.3 and 0.28);
   \node[font=\scriptsize] at (2.2,0.50) {$\mathcal{Y}(\Ldual w_\alpha)$};
   \fill (2.2,0.78) circle (2.2pt);
   \node[font=\scriptsize,anchor=south] at (2.2,0.84) {$p$};
   \node[font=\scriptsize] at (2.2,1.42) {$Z(\Ldual w)$：内部光滑};
   \draw[->,kwgray,dashed] (4.25,0.50) -- (4.75,0.50);
   \node[font=\footnotesize,anchor=west,align=left] at (4.80,0.50)
     {$\Ldual w_\alpha$ associated\\到 $\Ldual w$};
 \end{scope}
\end{tikzpicture}
\caption{单极子冒泡。$R^4$ 的尺度不变性让瞬子可以缩成一点；(a) 一个量子化单位的荷收缩到
't Hooft{} 算子位置并消失。(b) 在模空间上，这个过程把光滑内部与若干 associated 权型的
边界层连起来。紧化就是把这些边界层补上，于是修改空间=带奇异的 Schubert 簇。}
\label{fig:bogomolny}
\end{figure}

\paragraph{超凯勒结构出现在奇异点。}冒泡是局部现象，局部模型是 $C^2$ 上 $F$-等变瞬子模空间，
而 $R^4$ 上瞬子模空间自带超凯勒结构，$F$ 保持它，故等变瞬子模空间也是超凯勒的。
所以修改空间的奇异点附近有超凯勒结构（对 $SU(2)$ 情形就是 $A_1$ 奇异 $C^2/\ZZ_2$），
尽管修改空间整体不是超凯勒。\pA\anchor{src p.159}

\begin{kcard}{排除错误直觉 7：冒泡不是“微扰修正”}
\anchor{src p.156；src p.159}
\textbf{错误理解：}单极子冒泡是 't Hooft{} 算子边界条件的“非微扰修正”
（类似于瞬子修正），可以用 $e^{-S}$ 的幂次逐项估计。

\textbf{为什么错：}冒泡根本不是“对某个经典解的修正”，而是模空间自身的边界。
它出现在路径积分里，是因为路径积分遍历\emph{所有}场构型——包括那些把荷缩成点的构型。
它不是耦合展开的一项，没有“阶数”可数，也不是小量。原文明确指出：光滑单极子
（在 $R^3$ 上带通常边界条件）不会出现冒泡，冒泡\emph{只在有 't Hooft{} 算子时}发生。
\pA\anchor{src p.156}

\textbf{正确图像：}冒泡是“边界条件的非微扰修正”一语，应理解为：边界条件在经典层面
只规定了一个 GNO 权（一个奇异性），但在量子层面（即对整个模空间积分时），
边界附近还包含荷缩成点的拓扑扇区。这些扇区改变了“实际被积分到的”修改空间：
它从非紧的 $\mathcal{Y}(\Ldual w)$ 变成紧化的 Schubert 簇
$\overline{\mathcal{Y}}(\Ldual w)$。物理上，这正是 't Hooft{} 算子 OPE 中
次领头项的来源——它们全部来自奇异点。
\end{kcard}

\paragraph{自检问题。}
(1) 为什么 $F$-不变瞬子的冒泡点必须是 $C^2$ 的原点，而不是别的点？
(2) 若冒泡可用一维 $H_1$ 实现，为什么不能“只冒一半”（即 $s=1/2$ 的移动）？
(3) 用 \eqref{eq:pauly} 验证：当 $s=1$ 的冒泡把 $(m+1,m-1)$ 变成 $(m,m)$ 时，
修改空间维数是否改变？（提示：$\Delta$ 相同。）


\subsection{第六层：Kähler 商与 OPE 代数（§10.3--§10.4）}
\label{sec:flesh-hecke-kahler}

§10.1--§10.2 给了模空间 $Z$ 与它的边界。现在要回答第一个问题：$Z$ 与修改空间
$\mathcal{Y}$ 为什么\emph{相等}，而不只是维数相同？答案的结构非常漂亮：
\textbf{把 Bogomolny 方程重新解释为“全纯方程 $+$ 矩映射零点”}，
于是 $Z$ 是无穷维 Kähler 流形的辛商，而按复商/辛商对应它等于复商——
复商就是修改空间。

\subsubsection{§10.3：Bogomolny 方程 = 全纯方程 + 矩映射零点}

\paragraph{复结构。}在场空间 $\mathcal{C}=\{(A,\varphi_0)\}$ 上引入复结构 $I$，
取复坐标
\begin{equation}
A_{\bar z}\qquad\text{与}\qquad B=A_y-\ii\varphi_0.
\label{eq:cplxcoord}
\end{equation}
在这个复结构下，$F=\star D\varphi_0$ 的两个分量可以拆成
\begin{equation}
[D_y-\ii\varphi_0,\ D_{\bar z}]=0
\qquad\text{（全纯方程）},
\label{eq:holo}
\end{equation}
\begin{equation}
F_{23}=h\,D_y\varphi_0
\qquad\text{（实方程）}.
\label{eq:real}
\end{equation}
原文的说法是：(10.14) 在复结构 $I$ 下是全纯的。\pA\anchor{src p.160}

\paragraph{辛结构。}同一个场空间上放辛形式
\begin{equation}
\omega=\int_W\Tr\bigl(\delta A_2\wedge\delta A_3-h\,\delta A_1\wedge\delta\varphi_0\bigr)
\,dx^2dx^3dy,
\label{eq:omega}
\end{equation}
它与 $I$ 合成 Kähler 结构，度规为
$-\int_W\Tr(\delta A_2\otimes\delta A_2+\delta A_3\otimes\delta A_3
+h\delta A_y\otimes\delta A_y+h\delta\varphi_0\otimes\delta\varphi_0)$。
\pA\anchor{src p.160}

\paragraph{矩映射。}规范群在 $\mathcal{C}$ 上作用，由 $d\mu(v)=\iota_v\omega$ 定义矩映射。
对无穷小规范变换 $\delta A=-D\epsilon$、$\delta\varphi_0=[\epsilon,\varphi_0]$
（$\epsilon\in\Omega^0(C,\mathrm{ad}\,E)$），直接代入并分部积分：
\begin{align}
\iota_{\delta}\omega(\,\cdot\,)
&=\int_W\Tr\Bigl(\delta A_2\wedge\delta A_3-h\,\delta A_1\wedge\delta\varphi_0\Bigr)\notag\\
&=\int_W\Tr\Bigl((-D_2\epsilon)\,\delta A_3-(-D_3\epsilon)\,\delta A_2
-h(-D_1\epsilon)[\epsilon,\varphi_0]\Bigr)\notag\\
&=\int_W\Tr\Bigl(\epsilon\,(D_2\delta A_3-D_3\delta A_2-hD_y\delta\varphi_0)\Bigr)
\notag\\
&\qquad+\int_{\partial W}\Tr(\epsilon\,\varphi_0)\,h\,dx^2dx^3 .
\label{eq:moment}
\end{align}
（把 $D_2\delta A_3-D_3\delta A_2$ 组合成 $\delta F_{23}$，再用 $F_{23}=hD_y\varphi_0$ 的变分，
即得第一项。）故矩映射是
\begin{equation}
\mu(\epsilon)=\int_W\Tr\bigl(\epsilon(F_{23}-hD_y\varphi_0)\bigr)\,dx^2dx^3dy
+\int_{\partial W}\Tr(\epsilon\varphi_0)\,h\,dx^2dx^3 .
\label{eq:momentmap}
\end{equation}
\pA\anchor{src p.161}

\paragraph{关键一步：边界条件是被逼出来的。}要让 $\mu=0$ 恰好等于实 Bogomolny 方程
\eqref{eq:real}，边界项必须消失。而\eqref{eq:momentmap} 的边界项 $\Tr(\epsilon\varphi_0)$
有两个、也只有两个简单的消失方式：
\begin{itemize}
\item 令 $\epsilon=0$——即只允许在该边界上平凡的规范变换；
\item 令 $\varphi_0=0$——即在该边界上把 Higgs 场定为零。
\end{itemize}
§10.1 的边界条件正是“一端用一种、另一端用另一种”。原文因此说：
这些边界条件是“最简单的、允许把 Bogomolny 方程解释为全纯方程加矩映射零点的”
边界条件。\pA\anchor{src p.161}

\paragraph{Kähler 商 = 复商 = 修改空间。}于是
\begin{equation}
Z=\mu^{-1}(0)\big/\,G
\qquad\text{（辛商）},\qquad
Z'=\mathcal{C}_0\big/\,G^{\CC}
\qquad\text{（复商）},
\label{eq:quotients}
\end{equation}
其中 $\mathcal{C}_0$ 是满足全纯方程 \eqref{eq:holo} 的子空间。
通常比较辛商与复商要考虑不稳定/半稳定轨道，但这里因为规范变换在一端被要求平凡，
$G^{\CC}$ 自由作用且轨道闭，所有轨道都稳定，故 $Z=Z'$。\pA\anchor{src p.161}

最后识别 $Z'$：由 \eqref{eq:holo}，$E_y=E|_{\{y\}\times C}$ 的全纯型在 $y$ 不跨过
算子时不变，且这是无算子区间中唯一的 $G^{\CC}$-不变量；跨过算子时 $E_y$ 做 Hecke 修改。
故 $Z'$ 恰是 $E_-$ 的（指定型的）Hecke 修改空间。\pA\anchor{src p.162}

\begin{keyidea}
至此 M3 有了精确形式：
\begin{equation}
\underbrace{Z(\Ldual w;p)}_{\text{Bogomolny 解模空间}}
\;=\;\underbrace{\mu^{-1}(0)/G}_{\text{辛商}}
\;=\;\underbrace{\mathcal{C}_0/G^{\CC}}_{\text{复商}}
\;=\;\underbrace{\mathcal{Y}(\Ldual w;p)}_{\text{Hecke 修改空间}} .
\label{eq:m3exact}
\end{equation}
第一、二、三等号分别是：矩映射零点 $\Leftrightarrow$ 实方程；Kähler 商对应定理（无稳定性问题）；
全纯型沿 $y$ 不变。原文用一句话总结这件事的意义：这个复商/辛商比较
“为 't Hooft 与 Hecke 算子的关系提供了一个精确框架”。\pA\anchor{src p.162}
\end{keyidea}

\begin{kcard}{排除错误直觉 8：Kähler 商不是“形式地重写一遍”}
\anchor{src p.161；src p.162}
\textbf{错误理解：}既然 Bogomolny 方程本来就是“曲率 $=$ 某物”，把它写成矩映射零点
只是换个记号；商掉规范群也只是例行公事。

\textbf{为什么错：}矩映射零点不是一个恒等式，而是一个\emph{方程}：$\mu=0$ 恰好给出
$F_{23}=hD_y\varphi_0$——这是原来那个\emph{实}方程。而辛商的另一边，复商
$\mathcal{C}_0/G^{\CC}$ 是\emph{全纯}范畴的对象（全纯丛的模空间）。两者相等是
Kähler 几何的定理（在无稳定性问题时），不是定义。正是这个定理把“解微分方程”
变成“数全纯对象”，从而让 't Hooft{} 算子（微分方程侧）与 Hecke 算子（全纯侧）
可以比较。\pA\anchor{src p.162}

\textbf{进一步的收获：}边界条件不再是“为了算得出来而选的”，而是由“要让矩映射
零点等于实方程”唯一逼出来的。原文在 §10.3 明确把这当作 §10.1 边界条件的“新理由”。
\end{kcard}

\paragraph{两个独立的形变方向。}把 't Hooft{} 算子位置写成 $p_i=y_i\times z_i$，引入修改联络
$A'=A-\ii\varphi_0\,dy$。在复结构 $I$ 下 $A'_{\bar z}$ 与 $A'_y$ 全纯、$A'_z$ 由共轭确定；
保持纤维化 $W=I\times C\to C$ 的微分同胚群 $D_I$（固定 $z$、把 $y$ 映为
$\tilde y(y;z,\bar z)$，在边界上与 $y$ 一致）保持 $I$ 与全纯方程，故 $Z$ 作为\emph{复流形}
在 $D_I$ 下不变；而 $\omega$ 用到 $C$ 的度量、不是 $D_I$-不变的，故 $Z$ 的\emph{辛结构}会变。
于是在 $y$ 方向滑动算子（不让同 $z$ 的两算子交叉）改变辛结构而不改复结构；反之固定 $y_i$
变动 $z_i$ 相当于特殊地变动 $C$ 的复结构，改变复结构而不改辛结构。原文用这个观察重新解释了
§9.2 的全部紧化/去奇异现象（Fig.~17、18）。\pA\anchor{src p.163}\anchor{src p.164}

\subsubsection{§10.4：OPE 的系数是表示论的多重数}

\paragraph{Wilson 的类比。}对静态 Wilson 线，有
\begin{equation}
W(\Ldual R,p)\cdot W(\Ldual R',p')\ \longrightarrow\
\bigoplus_{\alpha\in\mathcal{R}}N_\alpha\otimes W(\Ldual R_\alpha,p'),
\label{eq:opeW}
\end{equation}
其中 $N_\alpha=\mathrm{Hom}_{\Ldual G}(\Ldual R\otimes\Ldual R',\Ldual R_\alpha)$
是\emph{向量空间}，不是数。为什么是向量空间？因为静态线算子不是作用在既有 Hilbert 空间上的
算子，而是“量子化定义的一部分”：它给理论加了一个外源荷，而荷的 Hilbert 空间是
表示空间。若取“时间”为圆，取迹，$N_\alpha$ 才被它的维数 $n_\alpha$ 替代
（这正是 Wilson 圈 OPE (6.27) 中系数为整数的原因）。\pA\anchor{src p.168}\anchor{src p.169}

\paragraph{'t Hooft{} 的 OPE。}S-对偶预言同一形式：
\begin{equation}
T(\Ldual R,p)\cdot T(\Ldual R',p')\ \longrightarrow\
\bigoplus_{\alpha\in\mathcal{R}}N_\alpha\otimes T(\Ldual R_\alpha,p').
\label{eq:opeT}
\end{equation}
而且即使不用 S-对偶，由于 $\Psi=0$ 处 't Hooft{} 算子是唯一的线算子，这样的展开
\emph{必须}成立，只是 $\mathcal{R}$ 与 $N_\alpha$ 待定。$\mathcal{R}$ 容易定（§6.4 的简单论证），
难的是 $N_\alpha$。\pA\anchor{src p.169}

\paragraph{膜的选择。}取两个 $(B,A,A)$ 型膜：$B$ 支撑在
$U=\{(\text{固定全纯型 }E_-,\ \phi\text{ 任意})\}$，$B'$ 支撑在 $U'=\{(\text{任意 }E,\ \phi=0)\}$。
$U$ 与 $U'$ 在 $M_H$ 中横截相交于唯一点 $(E_-,0)$，且都是复结构 $I$ 下的复拉格朗日子流形，
故为 $(B,A,A)$ 型。加上 't Hooft{} 算子后，超对称构型空间正是
$Z(\Ldual R_1,p_1;\dots;\Ldual R_k,p_k)$。\pA\anchor{src p.170}

\paragraph{态空间与代数结构。}无瞬子修正的 A-模型，其物理态空间是时间无关超对称构型的
上同调，即 $Z$ 的（相交）上同调：
\begin{equation}
\mathcal{H}(\Ldual R_1,\dots,\Ldual R_k)=H^*\bigl(Z(\Ldual R_1,\dots,\Ldual R_k)\bigr).
\label{eq:Hspace}
\end{equation}
$p_i$ 的位置无关紧要（TQFT 给出平坦联络，而 $W$ 中 $n$ 点组空间单连通，无单值性）。
$Z$ 拓扑上是简单乘积 $\prod_i Z(\Ldual R_i)$，于是
\begin{equation}
\mathcal{H}(\Ldual R_1,\dots,\Ldual R_k)=\bigotimes_{i=1}^k\mathcal{H}(\Ldual R_i),
\label{eq:Htensor}
\end{equation}
而 OPE 的同构 \eqref{eq:opeT} 归结为
\begin{equation}
\mathcal{H}(\Ldual R_1)\otimes\mathcal{H}(\Ldual R_2)
=\bigoplus_{\alpha\in\mathcal{R}}N_\alpha\otimes\mathcal{H}(\Ldual R_\alpha).
\label{eq:Hprod}
\end{equation}
另一方面，S-对偶预言 $N_\alpha$ 就是表示张量积的多重度空间：
\begin{equation}
\Ldual R_1\otimes\Ldual R_2
=\bigoplus_{\alpha\in\mathcal{R}}N_\alpha\otimes\Ldual R_\alpha .
\label{eq:tensor}
\end{equation}
把 \eqref{eq:Hprod} 与 \eqref{eq:tensor} 比较：$\mathcal{H}(\Ldual R)$ 与 $\Ldual R$ 服从\emph{同一个}张量代数。原文引 [127]：
满足同一张量代数且维数有限的族，其之间的同构唯一（模 $\Ldual G$ 共轭）——
故 $\mathcal{H}(\Ldual R)\cong\Ldual R$ 作为表示。\pA\anchor{src p.172}\anchor{src p.173}

\paragraph{幽灵数给出 \texorpdfstring{$\Ldual G$}{LG} 的作用。}A-模型的幽灵数
\begin{equation}
K=d-\dim_{\CC}X\qquad(d\ \text{为上同调类次数})
\label{eq:ghost}
\end{equation}
（费米共轭把 $d\leftrightarrow2n-d$、$K\to-K$）。$\Ldual G$ 在 $\Ldual R$ 上的作用
在 $\mathcal{H}$ 上对应为张量代数的一个导子，即 $\Ldual{\mathfrak g}$ 的一个元素。
用 $\Ldual G=SU(N)$ 的基本表示 $V$ 定标：$\mathcal{Y}(V)\cong\PP^{N-1}$，
$H^*(\PP^{N-1})$ 的次数为 $0,2,\dots,2N-2$，故
\begin{equation}
K\big|_V=\mathrm{diag}(N-1,\ N-3,\ \dots,\ -(N-1)).
\label{eq:eigen}
\end{equation}
$SU(2)$ 情形即 $K\sim\mathrm{diag}(1,-1)$。\pA\anchor{src p.173}\anchor{src p.174}

\paragraph{相交上同调。}$Z$ 一般奇异，用的上同调是相交上同调 $IH^*$：
粗略地说它是 $X$ 光滑部分的 $L^2$ 上同调（模标准猜想），且按 $K=d-\dim_{\CC}X$ 分级。
对支撑在子簇 $Y\subset X$ 上的复形 $IC_Y$，$H^*(X,IC_Y)$ 是 $Y$ 光滑部分的 $L^2$ 上同调
（次数移位 $-\dim_{\CC}Y$）。在 Hecke 修改空间的情形，相关的子簇是
$\mathcal{Y}(\Ldual w_\alpha)$（$\Ldual w_\alpha$ 为 associated 权）：
$IC(\Ldual w_\alpha)$ 支撑在“冒泡把权从 $\Ldual w$ 降低到 $\Ldual w_\alpha$ 的子空间”上。
\pA\anchor{src p.174}\anchor{src p.175}

\begin{example}[$SU(2)$：把 OPE 算出来]
取 $\Ldual G=SU(2)$、$\Ldual w=(1/2,-1/2)$（二维表示）。此时 $2\otimes2=3\oplus1$，即
\begin{equation}
T(\Ldual w)\cdot T(\Ldual w)=T(\Ldual w+\Ldual w)\oplus T(0),
\qquad \Ldual w+\Ldual w=(1,-1)\ (\text{伴随}).
\label{eq:ope31}
\end{equation}
几何上：同一 $z$、不同 $y<y'$ 处做两次修改，得 $\PP^1$ 上的 $\PP^1$ 丛，其上同调四维、
次数 $0,2,2,4$，$K$ 本征值 $-2,0,0,2$。令 $y\to y'$，出现 $A_1$ 奇异（局部 $R^4/\ZZ_2$），
其消解可局部模型化为 $T^*\PP^1$（带渐近 $R^4/\ZZ_2$ 的完备超凯勒度规）；
其上 $L^2$ 调和 2-形式空间是\emph{一维}的、$K=0$，对应平凡表示。于是三个调和形式
收敛到 $Z^*=\mathrm{WCP}^2(1,1,2)$ 光滑部分（伴随项），第四个收敛到奇异点的 delta 形式
（平凡项，系数为一维空间）。\pA\anchor{src p.176}\anchor{src p.177}
\end{example}

\paragraph{一般的 OPE 与数学表述。}对一般权，
\begin{equation}
T(\Ldual w)\cdot T(\Ldual w')
=T(\Ldual w+\Ldual w')\oplus\bigoplus_{\alpha\in\mathcal{R}'}
N_\alpha\otimes T(\Ldual w_\alpha),
\label{eq:ope32}
\end{equation}
其中 $\Ldual w_\alpha$ 取遍 $\Ldual w+\Ldual w'$ 的 associated 权。当 $y\to y'$ 时
模空间“吹落”（blowdown），对应有理映射
\begin{equation}
\pi:\ Z(\Ldual w,z,y;\Ldual w',z,y')\longrightarrow Z(\Ldual w+\Ldual w').
\label{eq:blowdown}
\end{equation}
在相交上同调上取推下：
\begin{equation}
\pi_*\bigl(IC(\Ldual w)\star IC(\Ldual w')\bigr)
=IC(\Ldual w+\Ldual w')\oplus\bigoplus_\alpha \widetilde N_\alpha\otimes IC(\Ldual w_\alpha).
\label{eq:push}
\end{equation}
领头项系数是平凡的（一维），因为 $\pi$ 在光滑部分是同构；其余系数是向量空间
$\widetilde N_\alpha$。Lusztig、Ginzburg、Mirković--Vilonen 的定理
\pA\anchor{src p.178}的核心就是证明 $\widetilde N_\alpha\cong N_\alpha$
（即 \eqref{eq:ope32} 中的多重数空间）。\pA\anchor{src p.178}\anchor{src p.179}

\begin{keyidea}
\textbf{M3 的代数形式。}把 \eqref{eq:Hprod} 与 \eqref{eq:tensor} 相比，得到：
\begin{center}
Hecke 算子（及其 OPE）构成的代数 $\;\cong\;$ $\Ldual G$ 的表示环。
\end{center}
这就是原文 §10.4 的结论，也是几何 Langlands 中“Hecke 特征层”的物理来源：
Hecke 算子作用在膜上，其“乘法”就是表示张量积。原文明确指出，这部分数学
“起源于 Lusztig 的工作”，经 Ginzburg 与 Mirković--Vilonen 精化为张量范畴的等价。
\pA\anchor{src p.169}\anchor{src p.172}
\end{keyidea}

\begin{kcard}{排除错误直觉 9：OPE 系数是“数”，或 Hecke 代数是“乘法算子”}
\anchor{src p.168；src p.178}
\textbf{错误理解 A：}OPE 的系数就是一些数（比如 $SU(2)$ 的 $2\otimes2=3\oplus1$ 中系数是 1），
所以 't Hooft{} 的 OPE 就是普通的表示论 Clebsch--Gordan 系数。

\textbf{为什么错：}对\emph{静态}线算子，系数是\emph{向量空间} $N_\alpha$（多重数空间），
只有当“时间”取成圆、取迹之后才退化为维数 $n_\alpha$。原文专门用 Wilson 的例子
说明这一点：闭 Wilson 圈是 holonomy 的迹，迹把向量空间换成它的维数；开（静态）Wilson 线
不是。\pA\anchor{src p.169}

\textbf{错误理解 B：}Hecke 算子代数就是“在函数空间上乘一个核”，所以它平凡。

\textbf{为什么错：}Hecke 算子的“乘法”是 OPE，而 OPE 的次领头项\emph{全部}来自
模空间的奇异点（单极子冒泡）。换句话说是非微扰的几何在决定代数结构；
\eqref{eq:push} 中 $\widetilde N_\alpha$ 的计算是原文称为“难点”的部分。
\end{kcard}

\paragraph{自检问题。}
(1) 在 \eqref{eq:ghost} 中，$K$ 的本征值在 $\PP^{N-1}$ 上是 $N-1,N-3,\dots$；
这与 $SU(N)$ 基本表示的权有何关系？
(2) 为什么 \eqref{eq:ope31} 的领先项系数是 $1$，而 \eqref{eq:ope32} 的次领头项
系数是向量空间？（提示：$\pi$ 在何处是同构。）
(3) 若把 $W$ 换成 $S^1\times C$（时间取圆），\eqref{eq:opeT} 的 $N_\alpha$ 会变成什么？


\subsection{第七层：放松时间反演，以及怎样诚实地读这一节（§10.5）}
\label{sec:flesh-hecke-honest}

\subsubsection{§10.5：扩展 Bogomolny 方程与 Higgs 丛的 Hecke 修改}

\paragraph{为什么必须放松时间反演。}前六层一直在用两个假设：解时间无关、解时间反演不变。
原文说时间无关性可以一般地证明，但\emph{时间反演不变性}是严重的限制：它只适用于
那些被选来保证 $\phi=0$ 的膜（§10.4 的 $U,U'$ 就是如此）。而几何 Langlands 中最重要的
A-膜——Hitchin 纤维化的\emph{一般纤维}、以及 §11 的典范余各向同性膜——支撑集并不局限在
$\phi=0$ 上。所以要研究 't Hooft{} 算子对这些膜的作用，必须放弃时间反演。\pA\anchor{src p.179}

\paragraph{方程。}把 Higgs 场写成 $\varphi=\varphi_0\,dx^0+\pi^*\widetilde\varphi$、
规范场写成 $A=A_0\,dx^0+\widetilde A$（$\pi$ 为投影 $M=R\times W\to W$），
$t=1$ 的时间无关 BPS 方程是
\begin{equation}
\widetilde F-\widetilde\varphi\wedge\widetilde\varphi
=\star\bigl(D\varphi_0-[A_0,\widetilde\varphi]\bigr),\qquad
\star D\widetilde\varphi=[\varphi_0,\widetilde\varphi]+DA_0,\qquad
D^*\widetilde\varphi+[A_0,\varphi_0]=0,
\label{eq:extbps}
\end{equation}
其中 $\star$、$D$、$D^*$ 均为三维意义。\pA\anchor{src p.179}

\paragraph{同一条消失定理。}若 $\varphi$ 在某个膜上满足 Dirichlet 或 Neumann，
则与 §\ref{sec:flesh-hecke-claim} 的 \eqref{eq:phising} 同一套积分论证给出
$\varphi_1=A_0=0$。取规范 $\widetilde A_y=0$、$W$ 上度规 $h|dz|^2+dy^2$，
\eqref{eq:extbps} 化为
\begin{equation}
D_{\bar z}\varphi_z=0,\qquad
\partial_yA_{\bar z}=-\ii D_{\bar z}\varphi_0,\qquad
\partial_y\varphi_z=-\ii[\varphi_z,\varphi_0],\qquad
F_{z\bar z}-[\varphi_z,\varphi_{\bar z}]=\frac{\ii h}{2}\partial_y\varphi_0 .
\label{eq:extred}
\end{equation}
\pA\anchor{src p.179}\anchor{src p.180}

\paragraph{第一个方程：Higgs 丛。}$D_{\bar z}\varphi_z=0$ 说 $\phi=\varphi_z\,dz$ 限制到
$C_y=\{y\}\times C$ 上时，是 $\mathrm{ad}(E)\otimes K_C$ 的全纯截面（对每个 $y$）。
故 $(E,\phi)$ 对每个 $y$ 都是 Higgs 丛。后两个方程说：$y$ 方向的演化只改变 $(E,\phi)$
的规范等价类，直到遇到 't Hooft{} 算子。\pA\anchor{src p.180}

\paragraph{关键区别：$\phi$ 不能被修改。}跨过算子时 $E$ 做通常的 Hecke 修改（同 §9.1）；
$(E,\phi)$ 也做“Hecke 修改”。但 $\phi$ 的角色与 $E$ 很不一样：
\begin{quote}
“全纯丛可以在一点被修改而保持全纯性，这正是 Hecke 或 't Hooft 算子所做的……
但一个全纯场，在这里是 $\mathrm{ad}(E)\otimes K_C$ 的全纯截面 $\phi$，不能在一点被修改
而保持全纯性。”
\pA\anchor{src p.180}
\end{quote}
所以给定 $z=0$ 之外的 $\phi$，它至多有一个全纯延拓。于是：加上 $\phi$ 不会给修改
\emph{增加}参数，反而一般会\emph{消掉}参数——因为 $y<y_0$ 时在 $z_0$ 全纯的 $\phi$，
在 $y>y_0$ 时一般会在 $z_0$ 出极点。

\paragraph{推导：参数怎么被消掉。}取 $G=U(N)$、权 $\Ldual w=(m_1,\dots,m_N)$。
$E_+$ 由全纯截面
\begin{equation}
t_i=(z-z_0)^{-m_i}s_i\qquad(i=1,\dots,N)
\label{eq:tmod}
\end{equation}
生成。设 $y<y_0$ 时 $\phi$ 在 $z_0$ 全纯，则
\begin{equation}
\phi\cdot s_i=\sum_jf_{ij}\,s_j\qquad(f_{ij}\ \text{在}\ z_0\ \text{全纯}),
\label{eq:phiS}
\end{equation}
于是由 \eqref{eq:tmod}
\begin{equation}
\phi\cdot t_i=\sum_j(z-z_0)^{m_j-m_i}f_{ij}\,t_j .
\label{eq:phiT}
\end{equation}
要让 $\phi$ 在 $y>y_0$ 时也全纯，需 $(z-z_0)^{m_j-m_i}f_{ij}$ 在 $z_0$ 正则。
排序 $m_1\ge\cdots\ge m_N$ 后，这对 $f_{ij}$ 的 Taylor 系数施加的条件数是
\begin{equation}
\sum_{i<j}(m_i-m_j),
\label{eq:phicount}
\end{equation}
而它恰好等于修改空间维数 \eqref{eq:dimformula}。所以修改空间的维度被\emph{全部消掉}：
\begin{center}
（$E$ 的修改空间维数）$-$（$\phi$ 全纯性的条件数）$=\ \Delta-\Delta=0$。
\end{center}
原文由此说：Higgs 丛 $(E,\phi)$ 的 Hecke 修改“一般是一个有限集”。
若 $\phi(z_0)$ 正则半单（$U(N)$ 情形即对角化且特征值互异），确实如此。
最系统的说法是：$\phi$ 在修改空间 $\mathcal{Y}(\Ldual w,z_0)$ 上决定一个向量场，
其零点就是可能的修改。（向量场不完全典范，但零点集是。）\pA\anchor{src p.181}

\begin{example}[$U(2)$，$\Ldual w=(1,0)$：数出两个点]
$\mathcal{Y}((1,0))\cong\PP^1$，一般截面
$s_+=(z-z_0)^{-1}u+(h_1,h_2)$，$u$ 模 $\CC^*$ 参数化 $\PP^1$。
$\phi$ 在 $y>y_0$ 全纯 $\iff$ $\phi\cdot s_+$ 仍具此形状 $\iff$ $\phi u\parallel u$。
正则半单时 $u$ 必须是 $\phi$ 的两个本征向量之一：修改空间从 $\PP^1$ 塌成\emph{两个点}。
\pA\anchor{src p.182}
\end{example}

\paragraph{直接证明“磁特征膜”。}由
$\partial_y\varphi_z=-\ii[\varphi_z,\varphi_0]$ 立即得到：若 $\varphi_z$ 在某个 $y$ 为零，
则处处为零（这正是 §10.4 用过的假设）。更重要的推论是：$\phi$ 的\emph{特征多项式}
与 $y$ 无关，甚至跨过 't Hooft{} 算子也不变。因此
\begin{center}
\textbf{Higgs 丛的 Hecke 变换把 Hitchin 纤维化的每根纤维映到自身。}
\end{center}
于是：支撑在 Hitchin 纤维上、带平坦 Chan--Paton 丛的膜，其支撑集在 Hecke 变换下不变。
按 S-对偶它必须是磁特征膜；而现在我们\emph{不需要} S-对偶，就得到了“支撑集不变”
这一半结论（完整证明见 [26]）。这也是本轮讨论班最值得注意的收获之一：
几何 Langlands 的“特征”性质，从扩展 Bogomolny 方程里直接看出来了。\pA\anchor{src p.182}

\paragraph{Hecke 对应。}与 §10.1 对称地，可以描述 Higgs 丛的 Hecke 对应：
在 $W=I\times C$ 两端都指定 $(A,\phi)$（\emph{不}指定 $\varphi_0$），除以两端平凡的
规范变换，解扩展 Bogomolny 方程。令 $W_0=\{(A,\phi):\ D\phi=0\}$（$M_H$ 的规范论版本），
则解得 $Q\subset W_0\times W_0$。$Q$ 一般很复杂，但“好区域”（每个 Higgs 丛只有
有限个该型的 Hecke 修改）中，$Q$ 是 $W_0\times W_0$ 的子流形，除规范群后是
$M_H\times M_H$ 的中维子流形。\pA\anchor{src p.182}\anchor{src p.183}

$(B,A,A)$ 型要求 $Q$ 在复结构 $I$ 下全纯，且对 $J,K$ 拉格朗日——等价于对全纯辛形式
$\Omega_I=\omega_J+\ii\omega_K$ 拉格朗日。更精确地，$Q$ 必须对
$\Omega=\pi_1^*\Omega_I-\pi_2^*\Omega_I$ 拉格朗日（负号使对角线上
“无算子的平凡 Hecke 对应”成为拉格朗日）。限制到 $Q$ 上，
\begin{equation}
\frac{1}{\pi}\Bigl(\int_{C_+}-\int_{C_-}\Bigr)|d^2z|\,\Tr(\delta\varphi_z\wedge\delta A_{\bar z})
=\frac{1}{\pi}\int_I dy\int_C|d^2z|\,\partial_y\Tr(\delta\varphi_z\wedge\delta A_{\bar z}).
\label{eq:Qform}
\end{equation}
用 \eqref{eq:extred} 的线性化
\begin{equation}
D_{\bar z}\delta\varphi_z=-[\delta A_{\bar z},\varphi_z],\qquad
\partial_y\delta A_{\bar z}=-\ii D_{\bar z}\delta\varphi_0-\ii[\delta A_{\bar z},\varphi_0],\qquad
\partial_y\delta\varphi_z=-\ii[\delta\varphi_z,\varphi_0]-\ii[\varphi_z,\delta\varphi_0],
\label{eq:Qlin}
\end{equation}
以及 Jacobi 恒等式与分部积分，可验证 \eqref{eq:Qform} 为零。故 $Q$ 是复拉格朗日的。
\pA\anchor{src p.183}

\subsubsection{怎样诚实地读这一节}

到这里，主链 S7 的完整链条是：§6.2 定义算子（GNO 权）$\to$ §9.1 丛的全纯型跳跃
$\to$ §9.2--9.3 修改空间 $=$ Schubert 胞腔 $\to$ §10.1--10.3 $\mathcal{Z}=$ 修改空间
$\to$ §10.2 边界 $=$ 单极子冒泡 $\to$ §10.4 OPE $=$ 表示环 $\to$ §10.5 扩展 Bogomolny
与 Hitchin 纤维不变性。但必须说清楚：这条链条上\emph{哪些环节是定理，哪些是物理论证}。
这是全文“物理论证代替严格证明”最严重的地方，我们判断（\pC）读者应按下表对待：

\begin{table}[htbp]
\centering\small
\begin{tabular}{@{}p{4.9cm}p{2.0cm}p{6.2cm}@{}}
\toprule
断言 & 强度 & 依据 \\
\midrule
't Hooft{} 算子由 $\Ldual G$ 的表示分类 & 定理 & GNO 对偶的算子版本（原文 \pA）\\
修改空间 $=$ $G(\mathcal{O})$-轨道 $=$ Schubert 胞腔 & 定理 & [122]，原文 \pA \\
维数公式 \eqref{eq:dimformula} & 定理 & 极点计数（本节补全）\\
指标公式 \eqref{eq:pauly} & 定理 & Pauly 用等变指标定理（原文 \pA）\\
冒泡 $=$ associated 权移动 & 定理 & ADHM $+$ 原文自己的计算（\pA）\\
$Z=$ 修改空间 & 半严格 & Kähler 商 $+$ 消失定理；无穷维细节见 [26]（\pB）\\
OPE 系数 $=$ 张量积多重度空间 & 定理 & Lusztig--Ginzburg--Mirković--Vilonen；原文只作物理翻译（\pB）\\
Higgs 丛修改有限性 & 半严格 & $\phi$ 正则半单假设 $+$ 向量场零点（\pB）\\
$Q$ 复拉格朗日 & 半严格 & 积分恒等式，边界项未逐一验证（\pB）\\
\bottomrule
\end{tabular}
\caption{本节各断言的强度。“定理”指数学上有完整证明（无论原文是否引用），
“半严格”指论证链完整但包含未经逐条验证的分析性假设。}
\label{tab:honest}
\end{table}

\paragraph{我们判断（\pC）数学读者应当这样读。}
把 §9--§10 当作“词典 $+$ 预言生成器”，而不是证明。它的价值在于：
(a) 把四个物理对象（'t Hooft{} 算子的奇异性、GNO 权、仿射 Grassmann 的点、表示环）
显式地写成同一个标号的四种化身；(b) 给出可独立检验的定量预言（维数公式、
OPE 系数、Hitchin 纤维不变性）；(c) 指出严格数学应该证什么（即
$\widetilde N_\alpha\cong N_\alpha$，这正是 [96--98] 的内容）。
反过来，若把“物理推导”当作证明，就会在“$Z=Y$”这类等式上过度自信：
那里的无穷维辛商、模空间紧化、相交上同调的使用都是\emph{形式}层面的，
原文自己也把细节推到 [26]。

\paragraph{最薄弱的两处。}\label{par:weakest}
\begin{enumerate}
\item \textbf{无穷维辛商与 $Z=Y$。}原文用“因规范变换限于一端平凡，稳定性问题被避开”
一句带过，但这句话针对的是有限维 Kähler 商的标准定理；这里空间无穷维、非紧、且带奇异，
“所有轨道稳定”需要证明。我们判断这是本节最脆弱的一环（\pC）。
\item \textbf{扩展 Bogomolny 的消失定理与有限性。}\eqref{eq:extbps}$\to$\eqref{eq:extred}
用到 $\varphi_1=A_0=0$ 的消失定理，而该定理依赖“膜上 Dirichlet 或 Neumann”
与“远处衰减”的假设；对一般膜这些假设并未验证。Higgs 丛修改的有限性也依赖
$\phi(z_0)$ 正则半单，这一条件不在一般 Langlands 情形中成立。\pC
\end{enumerate}

\begin{keyidea}
\textbf{三句话总结本节。}
(1) 't Hooft{} 算子是“规定绕线磁荷的奇异算符”，其 GNO 权 $\Ldual w$ 就是 Hecke 修改的型；
(2) 几何上，跨过奇异性时规范丛在一点被“局部重粘”，修改空间是仿射 Grassmann 的
Schubert 胞腔，其边界由单极子冒泡给出，其代数由 OPE 给出 $\Ldual G$ 的表示环；
(3) 物理证明的硬核是指标理论、ADHM 与仿射 Grassmann，软肋是无穷维辛商与消失定理——
这正是下一节（S8，D-模与特征层）要补的地方。
\end{keyidea}

\paragraph{本节总自检。}
(1) 若把 $G$ 换成 $P\,SU(2)$，\eqref{eq:dimformula} 与 \eqref{eq:pauly} 该怎样改？
（提示：$m_i$ 取半整数。）
(2) 为什么“$\phi$ 不能在一点被修改”这条观察，恰好把修改空间的维数消到零，
而不是消到某个正维数？
(3) 若 $\phi(z_0)$ 不正则半单（例如有重特征值），Higgs 丛的 Hecke 修改空间会变成什么？
它与 §10.2 的冒泡有什么关系？
